% Options for packages loaded elsewhere
\PassOptionsToPackage{unicode}{hyperref}
\PassOptionsToPackage{hyphens}{url}
\PassOptionsToPackage{dvipsnames,svgnames,x11names}{xcolor}
%
\documentclass[
  11pt,
]{article}
\usepackage{amsmath,amssymb}
\usepackage{iftex}
\ifPDFTeX
  \usepackage[T1]{fontenc}
  \usepackage[utf8]{inputenc}
  \usepackage{textcomp} % provide euro and other symbols
\else % if luatex or xetex
  \usepackage{unicode-math} % this also loads fontspec
  \defaultfontfeatures{Scale=MatchLowercase}
  \defaultfontfeatures[\rmfamily]{Ligatures=TeX,Scale=1}
\fi
\usepackage{lmodern}
\ifPDFTeX\else
  % xetex/luatex font selection
\fi
% Use upquote if available, for straight quotes in verbatim environments
\IfFileExists{upquote.sty}{\usepackage{upquote}}{}
\IfFileExists{microtype.sty}{% use microtype if available
  \usepackage[]{microtype}
  \UseMicrotypeSet[protrusion]{basicmath} % disable protrusion for tt fonts
}{}
\makeatletter
\@ifundefined{KOMAClassName}{% if non-KOMA class
  \IfFileExists{parskip.sty}{%
    \usepackage{parskip}
  }{% else
    \setlength{\parindent}{0pt}
    \setlength{\parskip}{6pt plus 2pt minus 1pt}}
}{% if KOMA class
  \KOMAoptions{parskip=half}}
\makeatother
\usepackage{xcolor}
\usepackage[margin=25mm]{geometry}
\usepackage{color}
\usepackage{fancyvrb}
\newcommand{\VerbBar}{|}
\newcommand{\VERB}{\Verb[commandchars=\\\{\}]}
\DefineVerbatimEnvironment{Highlighting}{Verbatim}{commandchars=\\\{\}}
% Add ',fontsize=\small' for more characters per line
\newenvironment{Shaded}{}{}
\newcommand{\AlertTok}[1]{\textcolor[rgb]{1.00,0.00,0.00}{\textbf{#1}}}
\newcommand{\AnnotationTok}[1]{\textcolor[rgb]{0.38,0.63,0.69}{\textbf{\textit{#1}}}}
\newcommand{\AttributeTok}[1]{\textcolor[rgb]{0.49,0.56,0.16}{#1}}
\newcommand{\BaseNTok}[1]{\textcolor[rgb]{0.25,0.63,0.44}{#1}}
\newcommand{\BuiltInTok}[1]{\textcolor[rgb]{0.00,0.50,0.00}{#1}}
\newcommand{\CharTok}[1]{\textcolor[rgb]{0.25,0.44,0.63}{#1}}
\newcommand{\CommentTok}[1]{\textcolor[rgb]{0.38,0.63,0.69}{\textit{#1}}}
\newcommand{\CommentVarTok}[1]{\textcolor[rgb]{0.38,0.63,0.69}{\textbf{\textit{#1}}}}
\newcommand{\ConstantTok}[1]{\textcolor[rgb]{0.53,0.00,0.00}{#1}}
\newcommand{\ControlFlowTok}[1]{\textcolor[rgb]{0.00,0.44,0.13}{\textbf{#1}}}
\newcommand{\DataTypeTok}[1]{\textcolor[rgb]{0.56,0.13,0.00}{#1}}
\newcommand{\DecValTok}[1]{\textcolor[rgb]{0.25,0.63,0.44}{#1}}
\newcommand{\DocumentationTok}[1]{\textcolor[rgb]{0.73,0.13,0.13}{\textit{#1}}}
\newcommand{\ErrorTok}[1]{\textcolor[rgb]{1.00,0.00,0.00}{\textbf{#1}}}
\newcommand{\ExtensionTok}[1]{#1}
\newcommand{\FloatTok}[1]{\textcolor[rgb]{0.25,0.63,0.44}{#1}}
\newcommand{\FunctionTok}[1]{\textcolor[rgb]{0.02,0.16,0.49}{#1}}
\newcommand{\ImportTok}[1]{\textcolor[rgb]{0.00,0.50,0.00}{\textbf{#1}}}
\newcommand{\InformationTok}[1]{\textcolor[rgb]{0.38,0.63,0.69}{\textbf{\textit{#1}}}}
\newcommand{\KeywordTok}[1]{\textcolor[rgb]{0.00,0.44,0.13}{\textbf{#1}}}
\newcommand{\NormalTok}[1]{#1}
\newcommand{\OperatorTok}[1]{\textcolor[rgb]{0.40,0.40,0.40}{#1}}
\newcommand{\OtherTok}[1]{\textcolor[rgb]{0.00,0.44,0.13}{#1}}
\newcommand{\PreprocessorTok}[1]{\textcolor[rgb]{0.74,0.48,0.00}{#1}}
\newcommand{\RegionMarkerTok}[1]{#1}
\newcommand{\SpecialCharTok}[1]{\textcolor[rgb]{0.25,0.44,0.63}{#1}}
\newcommand{\SpecialStringTok}[1]{\textcolor[rgb]{0.73,0.40,0.53}{#1}}
\newcommand{\StringTok}[1]{\textcolor[rgb]{0.25,0.44,0.63}{#1}}
\newcommand{\VariableTok}[1]{\textcolor[rgb]{0.10,0.09,0.49}{#1}}
\newcommand{\VerbatimStringTok}[1]{\textcolor[rgb]{0.25,0.44,0.63}{#1}}
\newcommand{\WarningTok}[1]{\textcolor[rgb]{0.38,0.63,0.69}{\textbf{\textit{#1}}}}
\usepackage{longtable,booktabs,array}
\usepackage{calc} % for calculating minipage widths
% Correct order of tables after \paragraph or \subparagraph
\usepackage{etoolbox}
\makeatletter
\patchcmd\longtable{\par}{\if@noskipsec\mbox{}\fi\par}{}{}
\makeatother
% Allow footnotes in longtable head/foot
\IfFileExists{footnotehyper.sty}{\usepackage{footnotehyper}}{\usepackage{footnote}}
\makesavenoteenv{longtable}
\setlength{\emergencystretch}{3em} % prevent overfull lines
\providecommand{\tightlist}{%
  \setlength{\itemsep}{0pt}\setlength{\parskip}{0pt}}
\setcounter{secnumdepth}{5}
\ifLuaTeX
  \usepackage{selnolig}  % disable illegal ligatures
\fi
\usepackage{bookmark}
\IfFileExists{xurl.sty}{\usepackage{xurl}}{} % add URL line breaks if available
\urlstyle{same}
\hypersetup{
  pdftitle={Concrete Lebesgue and Hausdorff Integral Applications},
  colorlinks=true,
  linkcolor={blue},
  filecolor={Maroon},
  citecolor={Blue},
  urlcolor={blue},
  pdfcreator={LaTeX via pandoc}}

\title{Concrete Lebesgue and Hausdorff Integral Applications}
\author{}
\date{15 October 2026}

\begin{document}
\maketitle

{
\hypersetup{linkcolor=}
\setcounter{tocdepth}{3}
\tableofcontents
}
\section{Concrete Lebesgue and Hausdorff integral
applications}\label{start}

A first collection of \textbf{16 fully worked examples}, written for an
analysis student. Start with the \hyperref[problems]{problem index},
then open the separate solution that you want. The examples favor
explicit computations, stated hypotheses, and direct arguments
explaining why the integration theorem applies.

\begin{itemize}
\tightlist
\item
  \hyperref[conventions]{Measures and theorem conventions}: normalized
  Hausdorff measure, completed-product sections, signed integrability,
  area and coarea.
\item
  \hyperref[problems]{Problem index}: no computed answers.
\item
  \hyperref[atlas]{Counterexample atlas}: precise failures and
  sufficient conditions for exchanging sums, limits, derivatives, and
  integrals.
\item
  \href{https://github.com/Drsunflowerherbert/Codex-for-Math/tree/codex/integral-examples-2026-10-15/integral-examples/pdf/integral-examples.pdf}{Full
  worked collection, PDF} and
  \href{https://github.com/Drsunflowerherbert/Codex-for-Math/tree/codex/integral-examples-2026-10-15/integral-examples/pdf/integral-examples.tex}{editable
  LaTeX}.
\item
  \href{https://github.com/Drsunflowerherbert/Codex-for-Math/tree/codex/integral-examples-2026-10-15/integral-examples/pdf/guide.pdf}{Compact
  guide, PDF} and
  \href{https://github.com/Drsunflowerherbert/Codex-for-Math/tree/codex/integral-examples-2026-10-15/integral-examples/pdf/guide.tex}{LaTeX}.
\item
  \hyperref[bibliography]{Source references},
  \href{https://github.com/Drsunflowerherbert/Codex-for-Math/tree/codex/integral-examples-2026-10-15/integral-examples/checks/source-references.json}{source
  snapshot hashes}, and \hyperref[verification]{verification ledger}.
\item
  \href{https://github.com/Drsunflowerherbert/Codex-for-Math/tree/codex/integral-examples-2026-10-15/integral-examples/index.json}{Machine-readable
  index} without results, and
  \href{https://github.com/Drsunflowerherbert/Codex-for-Math/tree/codex/integral-examples-2026-10-15/integral-examples/metadata.json}{full
  metadata} with hypotheses, results, and review status.
\end{itemize}

\subsection{Find a technique again}\label{find-a-technique-again}

From the repository root, using Python 3 with its standard library:

\begin{Shaded}
\begin{Highlighting}[]
\ExtensionTok{python3}\NormalTok{ integral{-}examples/find\_examples.py }\AttributeTok{{-}{-}topic}\NormalTok{ coarea}
\ExtensionTok{python3}\NormalTok{ integral{-}examples/find\_examples.py }\AttributeTok{{-}{-}technique}\NormalTok{ multiplicity}
\ExtensionTok{python3}\NormalTok{ integral{-}examples/find\_examples.py }\AttributeTok{{-}{-}prerequisite}\NormalTok{ dominated}
\ExtensionTok{python3}\NormalTok{ integral{-}examples/find\_examples.py }\AttributeTok{{-}{-}text}\NormalTok{ concentration}
\ExtensionTok{python3}\NormalTok{ integral{-}examples/find\_examples.py }\AttributeTok{{-}{-}list{-}tags}
\end{Highlighting}
\end{Shaded}

Filters are case-insensitive substring matches and combine with AND.
They print titles, tags, and paths relative to
\texttt{integral-\allowbreak{}examples/\allowbreak{}}. The \texttt{-\/-text} filter searches IDs,
titles, paths, topics, prerequisites, and techniques; it does not search
solution prose or recorded hypotheses/results. Add
\texttt{-\/-solutions} to print the recorded result. The files remain
ordinary Markdown and JSON; retrieval needs no hosted service.

\subsection{How to study with the
collection}\label{how-to-study-with-the-collection}

For each problem, identify the function, its domain, the measure, and
the integrability question before computing. In a region problem, write
the indicator and determine its sections. For signed integrals, examine
the absolute value or the positive and negative parts. For a
parametrization, compute the Jacobian and count preimages before
identifying the answer with an integral over the image. For a parameter
limit, write the proposed majorant explicitly.

E01--E06 cover Lebesgue integration and iterated integrals; E07--E13
cover curves, surfaces, area, coarea, and singular maps; E14--E16 cover
limit and derivative exchanges and approximate identities. Cross-links
in the atlas make failure cases easy to retrieve alongside their
positive criteria.

\subsection{Reproduce the checks and
PDFs}\label{reproduce-the-checks-and-pdfs}

\begin{Shaded}
\begin{Highlighting}[]
\ExtensionTok{python3}\NormalTok{ integral{-}examples/checks/check\_examples.py}
\ExtensionTok{python3}\NormalTok{ integral{-}examples/checks/check\_collection.py}
\ExtensionTok{python3}\NormalTok{ integral{-}examples/build.py}
\end{Highlighting}
\end{Shaded}

The concrete identity checker uses SymPy and mpmath. The collection
checker and retrieval command use only the Python standard library. The
build uses Pandoc and pdfLaTeX. See
\href{https://github.com/Drsunflowerherbert/Codex-for-Math/tree/codex/integral-examples-2026-10-15/integral-examples/checks/requirements.txt}{checks/requirements.txt}
for the versions actually used. The checked-in PDF and LaTeX need no
build tools to read.

Independent model reviewers audited every example; the ledger records
the scope, findings, and revisions. Symbolic and numerical checks
supplement those reviews and are \textbf{not proofs of analytic
hypotheses}. No Lean verification or human peer review is claimed.
General measure-theoretic theorems are cited prerequisites; this
collection does not reprove all of measure theory.

All collection changes are confined to \texttt{integral-\allowbreak{}examples/\allowbreak{}} on
branch \texttt{codex/\allowbreak{}integral-\allowbreak{}examples-\allowbreak{}2026-\allowbreak{}10-\allowbreak{}15}, based on the
integrated research commit
\texttt{e4f337c3\allowbreak{}aa52ec4c\allowbreak{}3f2aab07\allowbreak{}8d7c91b3\allowbreak{}2ee0015f}. Other research
branches are independent. The earlier paper was not used or
reconstructed. No third-party problem collection or source PDF is
reproduced, and no license is assigned here to the cited sources or to
unrelated repository work.

\section{Problem index --- without computed answers}\label{problems}

Each entry links to a separate full solution. Read
\hyperref[conventions]{CONVENTIONS.md} first for measure normalization
and the exact role of almost-everywhere sections. The JSON index and
retrieval command described in \hyperref[start]{README.md} support
search by topic, prerequisite, and technique without printing answers by
default.

\begin{longtable}[]{@{}
  >{\raggedright\arraybackslash}p{(\columnwidth - 6\tabcolsep) * \real{0.0700}}
  >{\raggedright\arraybackslash}p{(\columnwidth - 6\tabcolsep) * \real{0.3700}}
  >{\raggedright\arraybackslash}p{(\columnwidth - 6\tabcolsep) * \real{0.2500}}
  >{\raggedright\arraybackslash}p{(\columnwidth - 6\tabcolsep) * \real{0.3100}}@{}}
\toprule\noalign{}
\begin{minipage}[b]{\linewidth}\raggedright
ID
\end{minipage} & \begin{minipage}[b]{\linewidth}\raggedright
Task
\end{minipage} & \begin{minipage}[b]{\linewidth}\raggedright
Main prerequisites
\end{minipage} & \begin{minipage}[b]{\linewidth}\raggedright
Technique to look for
\end{minipage} \\
\midrule\noalign{}
\endhead
\bottomrule\noalign{}
\endlastfoot
\hyperref[E01]{E01} & Integrate \(e^{y^2}\) over \(0<x<y<1\). &
Measurability, Tonelli & Change order on a measurable region \\
\hyperref[E02]{E02} & Construct a region whose boundary has positive
area; integrate \(2y\) over it. & Countable additivity, Tonelli &
Explicit removed-interval sums \\
\hyperref[E03]{E03} & Can every section of a Lebesgue-measurable null
set be integrated? Prove the correct general statement. & Completion,
Fubini & Borel representative and null envelope \\
\hyperref[E04]{E04} & Compare the iterated integrals of
\((x^2-y^2)/(x^2+y^2)^2\) on \((0,1)^2\). & Signed integrals, Fubini &
Check positive and negative parts first \\
\hyperref[E05]{E05} & Sum an explicit signed matrix by rows and by
columns. & Counting measure, series & Detect infinite absolute mass \\
\hyperref[E06]{E06} & Evaluate and differentiate a Frullani parameter
integral. & Tonelli, dominated convergence & Replace a quotient by a
parameter integral \\
\hyperref[E07]{E07} & Integrate the first coordinate along a parabola
with respect to length. & Derivatives, area formula & Compute the speed
and pull back the weight \\
\hyperref[E08]{E08} & Integrate a coordinate square on a twice-traversed
circle. & Curve area formula & Count multiplicity; distinguish
circulation \\
\hyperref[E09]{E09} & Integrate height on a paraboloid graph over the
unit disk. & Gram determinant, polar coordinates & Derive the graph area
factor \\
\hyperref[E10]{E10} & Integrate a coordinate square over a sphere of
radius \(R\). & Surface area formula & Handle chart seams and pole
exceptions \\
\hyperref[E11]{E11} & Apply weighted area to a folded square
parametrization. & Jacobian, area formula & Compare source integral and
image multiplicity \\
\hyperref[E12]{E12} & Turn an integral on a disk into integrals over
circular level sets. & Gradient, coarea & Account for a critical point
before division \\
\hyperref[E13]{E13} & Explain why the Cantor staircase violates a
proposed derivative-area identity. & Cantor set, absolute continuity &
Locate the failed Lipschitz hypothesis \\
\hyperref[E14]{E14} & Track mass in concentration and escape examples
when taking limits under integrals. & Pointwise convergence, DCT &
Uniform integrability and spatial tightness \\
\hyperref[E15]{E15} & Compare a parameter derivative with the integral
of pointwise derivatives. & Difference quotients, DCT & Test local
domination; distinguish derivative limits \\
\hyperref[E16]{E16} & Recover a boundary value from exponential
averages. & Substitution, DCT & Approximate identity and an explicit
tail estimate \\
\end{longtable}

Suggested routes: E01--E06 for Lebesgue integration and iterated
integrals; E07--E13 for geometric integration; E14--E16 for exchanges of
limits, derivatives, and integrals. E03 and E04 are useful safeguards
before generalizing a successful calculation. The
\hyperref[atlas]{counterexample atlas} maps the safeguards to precise
sufficient conditions and contains some outcome spoilers.

\section{Measures, functions, and theorems used in this
library}\label{conventions}

The examples are applications for an analysis student. Their proofs are
written out; the general integration theorems below are prerequisites,
with sources in \hyperref[bibliography]{SOURCES.md}. A computation check
verifies an identity, not the measurability or integrability hypotheses
that justify it.

\subsection{Functions and values}\label{functions-and-values}

A function is a map with a specified domain, codomain, and rule. For
example, \(f:(0,1)^2\to\mathbb R\) is a function, \(f(x,y)\) is its
value at \((x,y)\), and \(f_x:y\mapsto f(x,y)\) is a section function.
With a fixed integrand and measure, the integral is a number (or an
allowed extended value), not another notation for the function. A
parameter integral defines a new function, for example
\(F:t\mapsto\int f_t\,d\mu\); \(F(t)\) is its value. When a section
integral is defined only almost everywhere, a measurable version is
chosen before integrating that version in the outer variable.

\subsection{Lebesgue and counting
measures}\label{lebesgue-and-counting-measures}

\(\mathcal L^n\) denotes complete \(n\)-dimensional Lebesgue measure on
\(\mathbb R^n\). A restricted integral \(\int_E f\,d\mathcal L^n\) means
integration over the measurable set \(E\) with the restricted measure.
For finite-valued \(f\) defined on all of \(\mathbb R^n\), it equals
\(\int f\mathbf1_E\,d\mathcal L^n\). Counting measure on
\(\mathbb N=\{1,2,\ldots\}\) turns integration into an explicit sum. The
Lebesgue and countable counting measure factors used in the product
theorems here are sigma-finite. This does not assert sigma-finiteness of
\(\mathcal H^k\) on all of \(\mathbb R^n\) when \(k<n\); our curve and
surface examples use finite-measure or countably covered restrictions.

For real measurable \(f\), set \(f^+=\max(f,0)\) and \(f^-=\max(-f,0)\).
A finite Lebesgue integral requires

\[
\int |f|\,d\mu=\int f^+\,d\mu+\int f^-\,d\mu<\infty.
\]

An extended signed integral is defined if at least one of the two terms
is finite. If both are infinite, their difference is undefined. An
improper limit or principal value must be named separately; its
existence does not supply a signed Lebesgue integral.

\subsection{Normalized Hausdorff
measure}\label{normalized-hausdorff-measure}

For an integer \(k\geq1\), let \(\omega_k=\mathcal L^k(B(0,1))\). We use

\[
\mathcal H^k(E)=\lim_{\delta\downarrow0}\inf\left\{
\sum_{i=1}^{\infty}\omega_k\left(\frac{\operatorname{diam}U_i}{2}\right)^k:
E\subseteq\bigcup_{i=1}^{\infty}U_i,\quad
\operatorname{diam}U_i\leq\delta\right\}.
\]

\(\mathcal H^0\) is counting measure. This normalization makes
\(\mathcal H^k\) agree with \(\mathcal L^k\) on every \(k\)-plane, using
an isometry to identify that plane with \(\mathbb R^k\). Thus
\(\mathcal H^1\) measures length and \(\mathcal H^2\) measures surface
area. In particular \(\omega_1=2\) and \(\omega_2=\pi\). Some books omit
the factor \(\omega_k/2^k\); their numerical values must be rescaled
before using these examples. All Hausdorff integrals in this library use
the convention in this paragraph.

\subsection{Tonelli and Fubini, including
completion}\label{tonelli-and-fubini-including-completion}

For a nonnegative product-measurable function on sigma-finite measure
spaces, Tonelli permits either order of integration, including the value
\(+\infty\). For a signed or complex product-measurable function, Fubini
applies if the absolute value has finite integral. It then gives finite
section integrals almost everywhere and equality of the two iterated
integrals with the product integral.

The product of two complete Lebesgue sigma-algebras need not be
complete. For \(\mathcal L^{p+q}\)-measurable sets or functions on
\(\mathbb R^{p+q}\), use the completed-product versions: sections need
be measurable only for almost every parameter. A Borel set has Borel
sections at every parameter; an arbitrary Lebesgue-measurable set need
not have measurable sections at every parameter. The outer integral uses
a measurable representative of the almost-everywhere section-integral
function. E03 proves this distinction directly with a null envelope. No
claim about all sections is inferred from completeness.

\subsection{Area with multiplicity}\label{area-with-multiplicity}

For a Lipschitz map \(F:\mathbb R^m\to\mathbb R^n\) with \(m\leq n\),
write

\[
J_mF(x)=\sqrt{\det(DF(x)^TDF(x))},\qquad
N(F,A,z)=\#\{x\in A:F(x)=z\}.
\]

The derivative and Jacobian exist almost everywhere; set the Jacobian to
zero at the exceptional points where the derivative does not exist. For
a measurable domain \(A\) and nonnegative Borel weight \(h\) on the
source, the weighted area formula is

\[
\int_A h(x)J_mF(x)\,d\mathcal L^m(x)
=\int_{\mathbb R^n}\left(\sum_{x\in A\cap F^{-1}(\{z\})}h(x)\right)
\,d\mathcal H^m(z).
\]

The sum is an extended nonnegative sum. With \(h=1\) it counts
multiplicity. With a nonnegative Borel \(w\) and \(h=w\circ F\) it
becomes \(w(z)N(F,A,z)\), with the nonnegative-integration convention
\(0\cdot\infty=0\). An injective parametrization therefore computes the
geometric scalar integral once. If multiplicity is \(q\) almost
everywhere, the parameter integral counts it \(q\) times. For signed
weights require finite absolute integral with this multiplicity. A map
on a bounded domain that is Lipschitz there can be extended to apply the
formula; the geometric parametrizations to which we apply area are also
Lipschitz on a suitable containing compact box. The Cantor staircase is
an explicitly excluded counterexample.

For a graph \(F(u)=(u,g(u))\) with a scalar \(C^1\) function \(g\),
\(J_mF=\sqrt{1+|\nabla g|^2}\). This factor is derived in the graph
example.

\subsection{Coarea and division by a
Jacobian}\label{coarea-and-division-by-a-jacobian}

For a Lipschitz map \(G:\mathbb R^n\to\mathbb R^m\), \(n\geq m\), the
coarea Jacobian is

\[
J_mG(x)=\sqrt{\det(DG(x)DG(x)^T)}.
\]

For nonnegative Borel \(h\), coarea says

\[
\int_{\mathbb R^n}h(x)J_mG(x)\,d\mathcal L^n(x)
=\int_{\mathbb R^m}\left(\int_{G^{-1}(\{t\})}h(x)
\,d\mathcal H^{n-m}(x)\right)\,d\mathcal L^m(t).
\]

Fiber statements and inner integrals have their usual almost-everywhere
meaning. Here \(m=1\) gives \(J_1G=|\nabla G|\). To recover an
unweighted volume integral, division by \(J_mG\) is legitimate only on
its positive set. Any contribution from its zero set must be handled
separately. E12 identifies and disposes of its critical point
explicitly. E13 shows why a singular continuous map cannot be
substituted into a Lipschitz area formula.

\subsection{Scalar integration and
orientation}\label{scalar-integration-and-orientation}

\(\int_S w\,d\mathcal H^k\) has no orientation. A scalar curve
parametrization uses speed \(|\gamma'(t)|\). A vector circulation uses
\(\int F(\gamma(t))\cdot\gamma'(t)\,dt\); reversing traversal changes
its sign. Flux through an oriented surface is
\(\int_S F\cdot\nu\,d\mathcal H^{n-1}\) with a chosen unit normal
\(\nu\). Changing \(\nu\) changes flux. Neither a tangent vector nor a
normal vector is part of Hausdorff measure itself.

\subsection{Convergence and
differentiation}\label{convergence-and-differentiation}

The version of monotone convergence used here concerns nonnegative
increasing measurable functions. Dominated convergence concerns
measurable functions and requires almost-everywhere convergence and one
integrable majorant for the entire sequence. Pointwise convergence,
boundedness of each individual function, or a uniform bound on the
integrals alone is insufficient. E14 gives concentration and escape
examples with explicit sufficient repairs. E15 states and proves a local
domination criterion for differentiating an integral. E16 explains how
concentration can instead be used deliberately as an approximate
identity.

\section{E01 --- Change the order to create the missing
derivative}\label{E01}

\subsection{Problem}\label{problem}

Let

\[
D=\{(x,y)\in\mathbb R^2:0<x<y<1\}.
\]

Compute \(\int_D e^{y^2}\,d\mathcal L^2(x,y)\) without finding an
antiderivative of \(y\mapsto e^{y^2}\).

\subsection{Data and hypotheses}\label{data-and-hypotheses}

The measure is two-dimensional Lebesgue measure \(\mathcal L^2\); the
one-dimensional integrals below use \(\mathcal L^1\). Define the
function on the whole plane by

\[
F(x,y)=\mathbf 1_D(x,y)e^{y^2}.
\]

Here \(F\) is a function, whereas \(F(x,y)\) is its value at the point
\((x,y)\). The region \(D\) is open and hence Borel measurable. The
function \(F\) is Borel measurable, is nonnegative, and satisfies
\(0\leq F\leq e\mathbf 1_{(0,1)^2}\). Consequently
\(\int_{\mathbb R^2}|F|\,d\mathcal L^2\leq e<\infty\) before any
interchange is made.

\subsection{Solution}\label{solution}

\textbf{1. Translate the region into sections.} Fix \(y\in(0,1)\). The
permitted values of \(x\) are exactly \(0<x<y\), so the horizontal
section has length \(y\). Fix \(x\in(0,1)\) instead. The permitted
values of \(y\) are \(x<y<1\). Thus the two descriptions are

\[
D=\{(x,y):0<y<1,\ 0<x<y\}
 =\{(x,y):0<x<1,\ x<y<1\}.
\]

These are descriptions of the same set, not a rule for swapping the
symbols \(dx\) and \(dy\) while keeping the bounds unchanged.

\textbf{2. Integrate in the useful direction.} Fubini's theorem,
justified by the absolute-integrability bound above, gives

\[
\begin{aligned}
\int_D e^{y^2}\,d\mathcal L^2(x,y)
&=\int_0^1\left(\int_0^y e^{y^2}\,dx\right)dy\\
&=\int_0^1 y e^{y^2}\,dy.
\end{aligned}
\]

In the inner integral, \(y\) is fixed, so \(e^{y^2}\) is constant as a
function of \(x\). Multiplying that value by the section length produces
the factor \(y\).

\textbf{3. Finish with a one-variable substitution.} Put \(u=y^2\), so
\(du=2y\,dy\). Then

\[
\int_0^1 y e^{y^2}\,dy
=\frac12\int_0^1e^u\,du
=\boxed{\frac{e-1}{2}}.
\]

Equivalently, this also evaluates the nested expression

\[
\int_0^1\left(\int_x^1 e^{y^2}\,dy\right)dx
=\frac{e-1}{2}.
\]

The inner integral in this last expression is perfectly well-defined;
lacking an elementary antiderivative does not make a function
nonintegrable.

\textbf{4. A sum check.} Uniform convergence of the exponential power
series on \([0,1]\) permits termwise integration here and gives

\[
\int_0^1 y e^{y^2}\,dy
=\sum_{k=0}^{\infty}\frac{1}{k!}\int_0^1y^{2k+1}\,dy
=\sum_{k=0}^{\infty}\frac{1}{2(k+1)k!}
=\frac12\sum_{j=1}^{\infty}\frac1{j!}.
\]

The last series is \((e-1)/2\). Alternatively, the nonnegative partial
sums and monotone convergence justify this same termwise calculation.

\subsection{Why the theorem applies}\label{why-the-theorem-applies}

Both factor measures are \(\sigma\)-finite. The region indicator makes
the varying domain into a measurable function on a fixed product space.
The bound \(|F|\leq e\mathbf 1_{(0,1)^2}\) verifies Fubini's
absolute-integrability hypothesis independently of the answer. Because
\(F\) is nonnegative, Tonelli also allows either order before finiteness
is known. Every section in this example is Borel measurable; the weaker
almost-everywhere assertion for arbitrary Lebesgue-measurable functions
is discussed in E03.

\subsection{Pitfalls}\label{pitfalls}

\begin{itemize}
\tightlist
\item
  Keeping \(x<y<1\) as the inner \(x\)-bounds is meaningless: the bounds
  must describe the variable being integrated.
\item
  The integral \(\int_D e^{y^2}\,d\mathcal L^2\) integrates a function
  over a set. It does not integrate the set or identify a function with
  one of its values.
\item
  Replacing strict inequalities by non-strict ones is harmless here
  because the added line segments have \(\mathcal L^2\)-measure zero.
  E02 explains why such a boundary argument needs verification for a
  general region.
\end{itemize}

\subsection{Reusable technique}\label{reusable-technique}

\begingroup\small\raggedright Tags: \texttt{region-\allowbreak{}sections}, \texttt{Fubini}, \texttt{Tonelli},
\texttt{indicator-\allowbreak{}extension}, \texttt{order-\allowbreak{}selection},
\texttt{power-\allowbreak{}series}.\par\endgroup

When an inner primitive looks unhelpful, write both section descriptions
explicitly. Integrating in the other direction can supply the derivative
factor needed for a substitution. First verify measurability and
nonnegativity or absolute integrability.

\subsection{Sources}\label{sources}

\begin{itemize}
\tightlist
\item
  John K. Hunter,
  \href{https://www.math.ucdavis.edu/~hunter/measure_theory/measure_notes.pdf}{\emph{Measure
  Theory}}, Theorem 5.18 (Fubini), printed p.~61; Theorem 4.6 (monotone
  convergence). Source ID: \texttt{hunter-\allowbreak{}measure}. The region,
  computation, and sum check above are worked out directly.
\end{itemize}

\section{E02 --- A measurable region whose entire boundary has positive
area}\label{E02}

\subsection{Problem}\label{problem-1}

Construct a closed set \(C\subset[0,1]\) with empty interior and
\(\mathcal L^1(C)=1/2\). For \(E=C\times[0,1]\), determine its boundary
and compute

\[
\int_E 2y\,d\mathcal L^2(x,y).
\]

Explain why simply discarding the boundary gives the wrong answer and
why the zero extension of this integrand is not Riemann integrable on
\([0,1]^2\).

\subsection{Data and hypotheses}\label{data-and-hypotheses-1}

All measures are Lebesgue measures \(\mathcal L^1\) and
\(\mathcal L^2\). Start with \(C_0=[0,1]\). At stage \(n\geq1\), remove
from the middle of each of the \(2^{n-1}\) closed intervals in
\(C_{n-1}\) one open interval of length \(4^{-n}\). Let \(C_n\) be the
remaining union of \(2^n\) closed intervals and put

\[
C=\bigcap_{n=0}^{\infty}C_n.
\]

This is the Smith--Volterra construction specified by these exact
lengths. The sets and integrands are explicit except for the infinite
intersection defining \(C\); no nonmeasurable choice is involved.

\subsection{Solution}\label{solution-1}

\textbf{1. Check that every removal fits.} All component intervals of
\(C_n\) have equal length

\[
\ell_n=2^{-n-1}+2^{-2n-1}.
\]

For \(n=0\) this is \(1\). If the formula holds at stage \(n-1\), the
next component length is

\[
\frac{\ell_{n-1}-4^{-n}}2
=2^{-n-1}+2^{-2n-1}>0.
\]

In particular the removed middle interval is strictly shorter than its
parent, so the construction is possible at every stage.

\textbf{2. Compute the measure by an explicit sum.} Removed intervals
from different stages are disjoint. Their total length through stage
\(n\) is

\[
\sum_{k=1}^{n}2^{k-1}4^{-k}
=\frac12\sum_{k=1}^{n}2^{-k}
=\frac12(1-2^{-n}).
\]

Therefore

\[
\mathcal L^1(C_n)=\frac12+2^{-n-1},\qquad
\mathcal L^1(C)=\lim_{n\to\infty}\mathcal L^1(C_n)=\frac12.
\]

The limit follows from continuity from above of a measure, since \(C_n\)
decreases and \(\mathcal L^1(C_0)=1<\infty\).

\textbf{3. Prove the interior is empty.} Suppose \(C\) contained a
nondegenerate interval \(J\). For every \(n\), the connected set \(J\)
would lie in one component interval of \(C_n\), because distinct
components are separated by a removed open interval. Hence the length of
\(J\) would be at most \(\ell_n\) for every \(n\). But \(\ell_n\to0\), a
contradiction. Since \(C\) is closed and has empty interior, its
complement is dense in \(\mathbb R\).

\textbf{4. Find the boundary of \(E\).} The set \(E=C\times[0,1]\) is
closed. It has empty interior: every open ball about a point of \(E\)
contains a point whose first coordinate lies outside \(C\). Thus

\[
\partial E=\overline E\setminus E^{\circ}=E,
\qquad \mathcal L^2(\partial E)=\mathcal L^2(E)=\frac12.
\]

The last equality uses the measurable-rectangle formula
\(\mathcal L^2(C\times[0,1])=\mathcal L^1(C)\mathcal L^1([0,1])\).

\textbf{5. Integrate.} Extend the integrand to the square as
\(F(x,y)=2y\mathbf 1_C(x)\). It is Borel measurable and
\(0\leq F\leq2\), so Fubini applies:

\[
\begin{aligned}
\int_E2y\,d\mathcal L^2(x,y)
&=\int_0^1\mathbf 1_C(x)\left(\int_0^1 2y\,dy\right)dx\\
&=\int_0^1\mathbf 1_C(x)\,dx
=\boxed{\frac12}.
\end{aligned}
\]

Deleting the boundary would leave the empty set
\(E\setminus\partial E\), whose integral is \(0\), not \(1/2\).

\textbf{6. Compare with Riemann integration directly.} Take any
rectangular partition of \([0,1]^2\). Every subrectangle of positive
width contains a point whose first coordinate is outside \(C\). Since
\(F\geq0\), its infimum on each such rectangle is \(0\). Every lower
Darboux sum is therefore \(0\). Every upper Darboux step function is at
least \(F\), so its integral, equal to the corresponding upper sum, is
at least \(\int F\,d\mathcal L^2=1/2\). Upper and lower sums cannot
approach a common value. Thus \(F\) is not Riemann integrable on the
square, although it is bounded and Lebesgue integrable.

\subsection{Why the theorem applies}\label{why-the-theorem-applies-1}

Closedness establishes Borel measurability. Continuity from above needs
the finite measure of the initial set; that hypothesis was checked. For
the integral, the bounded Borel function \(F\) on a finite-measure
square is absolutely integrable, so Fubini applies. No assumption that
the boundary is null is required by Lebesgue integration.

\subsection{Pitfalls}\label{pitfalls-1}

\begin{itemize}
\tightlist
\item
  ``Boundary'' does not imply ``measure zero.'' Here the entire
  positive-area set is its boundary.
\item
  Empty interior does not imply zero measure.
\item
  A bounded measurable function need not be Riemann integrable. The
  direct Darboux argument identifies the actual obstruction in this
  example.
\item
  The stage-\(n\) removal length \(4^{-n}\) is the length of
  \textbf{each} of \(2^{n-1}\) gaps, not their combined length.
\end{itemize}

\subsection{Reusable technique}\label{reusable-technique-1}

\begingroup\small\raggedright Tags: \texttt{fat-\allowbreak{}Cantor-\allowbreak{}set}, \texttt{positive-\allowbreak{}boundary-\allowbreak{}measure},
\texttt{explicit-\allowbreak{}series}, \texttt{continuity-\allowbreak{}from-\allowbreak{}above},
\texttt{Fubini}, \texttt{Riemann-\allowbreak{}versus-\allowbreak{}Lebesgue}.\par\endgroup

For a complicated closed region, calculate removed measure with a
disjoint sum, check interior separately, and integrate through an
indicator. Treat a proposed boundary deletion as a measure claim that
needs proof.

\subsection{Sources}\label{sources-1}

\begin{itemize}
\tightlist
\item
  John K. Hunter,
  \href{https://www.math.ucdavis.edu/~hunter/measure_theory/measure_notes.pdf}{\emph{Measure
  Theory}}, Proposition 1.12 (continuity of measures), Theorem 5.14
  (product measure), and Theorem 5.18 (Fubini). Source ID:
  \texttt{hunter-\allowbreak{}measure}. The construction, interval-length formula,
  boundary proof, and Darboux argument are supplied here.
\end{itemize}

\section{E03 --- A null measurable set with one nonmeasurable
section}\label{E03}

\subsection{Problem}\label{problem-2}

Find a Lebesgue-measurable set \(E\subset\mathbb R^2\) for which the
vertical section

\[
E_x=\{y\in\mathbb R:(x,y)\in E\}
\]

is not Lebesgue measurable for one value of \(x\). Reconcile this with
Fubini's theorem, and prove the correct almost-everywhere statement for
sections of arbitrary Lebesgue-measurable sets.

\subsection{Data and hypotheses}\label{data-and-hypotheses-2}

Use complete Lebesgue measures \(\mathcal L^1\) and \(\mathcal L^2\).
Choose a non-Lebesgue-measurable set \(V\subset[0,1]\) and define

\[
E=\{0\}\times V.
\]

Such a \(V\) exists under the usual axiom of choice. For completeness,
choose one representative from each equivalence class in \([0,1]\) under
\(s\sim t\) when \(s-t\in\mathbb Q\). The translates \(V+q\), for
\(q\in\mathbb Q\cap[-1,1]\), are disjoint, cover \([0,1]\), and lie in
\([-1,2]\). If \(V\) were measurable with measure zero, their countable
union could not cover \([0,1]\). If it had positive measure, countably
many disjoint translates would have infinite measure inside \([-1,2]\).
These two alternatives prove nonmeasurability.

\subsection{Solution}\label{solution-2}

\textbf{1. Joint measurability does hold.} The set \(E\) lies in the
Borel set \(N=\{0\}\times[0,1]\). For every \(\varepsilon>0\),

\[
N\subset(-\varepsilon,\varepsilon)\times[0,1],
\qquad \mathcal L^2(N)\leq2\varepsilon.
\]

Hence \(\mathcal L^2(N)=0\). Completeness of \(\mathcal L^2\) implies
that its subset \(E\) is Lebesgue measurable and has measure zero. Thus
the function \(\mathbf 1_E\) is Lebesgue measurable and belongs to
\(L^1(\mathbb R^2)\), with integral \(0\).

\textbf{2. Inspect each section.} At \(x=0\), the section is \(E_0=V\),
which is not Lebesgue measurable. For every \(x\ne0\),
\(E_x=\varnothing\), so its measure and the integral of its indicator
are \(0\). In particular

\[
y\longmapsto\mathbf 1_E(0,y)=\mathbf 1_V(y)
\]

is a nonmeasurable section function. One must not write its Lebesgue
integral as though it were defined.

For comparison, every horizontal section is either \(\{0\}\) or
\(\varnothing\), so it is measurable and has measure zero. In this
direction the section-measure function is identically \(0\), even though
deciding which of the two section sets occurs depends on membership in
\(V\).

\textbf{3. Interpret the iterated integral correctly.} On the
full-measure set \(\mathbb R\setminus\{0\}\) the vertical inner integral
exists and equals \(0\). Extend that outer integrand by the value \(0\)
at the exceptional point. Its integral is \(0\), agreeing with
\(\int\mathbf 1_E\,d\mathcal L^2=0\). This is an extension of an
almost-everywhere defined inner integral; it is not an assertion that
the original inner integral existed at \(x=0\).

\textbf{4. Prove the general section statement.} Let
\(A\subset\mathbb R^{m+n}\) be Lebesgue measurable. Borel regularity and
completion provide Borel sets \(B,N\subset\mathbb R^{m+n}\) such that

\[
A\mathbin{\triangle}B\subset N,
\qquad \mathcal L^{m+n}(N)=0.
\]

Here \(\triangle\) denotes symmetric difference. For every
\(x\in\mathbb R^m\), the sets \(B_x\) and \(N_x\) are Borel subsets of
\(\mathbb R^n\): the section is the preimage under the continuous map
\(y\mapsto(x,y)\). Since \(\mathbf 1_N\in L^1\), Fubini applied to this
Borel function gives

\[
\int_{\mathbb R^m}\mathcal L^n(N_x)\,d\mathcal L^m(x)=0.
\]

Its nonnegative integrand is therefore zero outside a measurable
\(\mathcal L^m\)-null set \(Z\). For \(x\notin Z\),

\[
A_x\mathbin{\triangle}B_x\subset N_x,
\qquad \mathcal L^n(N_x)=0.
\]

Completeness of \(\mathcal L^n\) now implies that \(A_x\) is Lebesgue
measurable and has the same measure as \(B_x\). Thus \textbf{vertical
sections are Lebesgue measurable for almost every \(x\)}. Interchanging
the coordinate blocks proves the corresponding horizontal assertion.
This argument does not prove measurability at parameters in \(Z\), and
the example shows why it cannot.

\subsection{Why the theorem applies}\label{why-the-theorem-applies-2}

In Step 4, Fubini is applied to the Borel null envelope \(N\), whose
indicator is absolutely integrable. It is not prematurely applied to a
possibly nonmeasurable section of \(A\). For an \(L^1\)
Lebesgue-measurable function, completed-measure Fubini similarly gives
measurable, integrable sections almost everywhere and equality after the
almost-everywhere defined outer integrand is extended on a null set.

A set in the uncompleted product \(\sigma\)-algebra
\(\mathcal L(\mathbb R^m)\otimes\mathcal L(\mathbb R^n)\) has measurable
sections at \textbf{every} parameter. Indeed, sets with measurable
sections form a \(\sigma\)-algebra containing all measurable rectangles.
The example \(E\) therefore belongs to the completed planar Lebesgue
\(\sigma\)-algebra but not to that uncompleted product
\(\sigma\)-algebra.

\subsection{Pitfalls}\label{pitfalls-2}

\begin{itemize}
\tightlist
\item
  ``Jointly Lebesgue measurable'' does not imply that every section is
  Lebesgue measurable.
\item
  Completeness of each factor measure does not make the uncompleted
  product measure complete.
\item
  Changing a jointly measurable function on a planar null set can damage
  exceptional one-dimensional sections. It does not change the joint
  Lebesgue integral.
\item
  This construction requires a nonmeasurable set; it is not an explicit
  numerical algorithm for testing membership in \(V\).
\end{itemize}

\subsection{Reusable technique}\label{reusable-technique-2}

\begingroup\small\raggedright Tags: \texttt{exceptional-\allowbreak{}sections}, \texttt{completion},
\texttt{Borel-\allowbreak{}representative}, \texttt{null-\allowbreak{}envelope}, \texttt{Fubini},
\texttt{Vitali-\allowbreak{}set}.\par\endgroup

To justify a section argument for completed measures, replace the joint
object by a Borel representative, control the difference inside a Borel
null envelope, and use Fubini on that envelope. Record the exceptional
parameter set rather than suppressing it.

\subsection{Sources}\label{sources-2}

\begin{itemize}
\tightlist
\item
  John K. Hunter,
  \href{https://www.math.ucdavis.edu/~hunter/measure_theory/measure_notes.pdf}{\emph{Measure
  Theory}}, Example 5.20 and Theorem 5.21, printed p.~62; Theorems
  2.27--2.28 (Borel regularity and completion). Source ID:
  \texttt{hunter-\allowbreak{}measure}. The null-set construction is the classical
  completion example; the representative-and-envelope proof and its
  integral interpretation are detailed above.
\end{itemize}

\section{E04 --- Two finite iterated integrals, but no signed Lebesgue
integral}\label{E04}

\subsection{Problem}\label{problem-3}

On \(Q=(0,1)^2\), define the real-valued function

\[
f(x,y)=\frac{x^2-y^2}{(x^2+y^2)^2}.
\]

Compute both iterated integrals. Determine whether
\(\int_Q f\,d\mathcal L^2\) exists as a signed Lebesgue integral, and
explain exactly which hypothesis prevents use of signed Fubini.

\subsection{Data and hypotheses}\label{data-and-hypotheses-3}

The measure on \(Q\) is the restriction of planar Lebesgue measure
\(\mathcal L^2\). Inner and outer one-variable integrals use
\(\mathcal L^1\). The function \(f\) is continuous on the open square,
so it is Borel measurable. Its singularity is approached as
\((x,y)\to(0,0)\), although that point is outside \(Q\).

For a real-valued measurable function, put \(f^+=\max(f,0)\) and
\(f^-=\max(-f,0)\). The signed Lebesgue integral is
\(\int f^+-\int f^-\) only when these two quantities are not both
infinite. A finite signed Lebesgue integral requires \(\int|f|<\infty\).

\subsection{Solution}\label{solution-3}

\textbf{1. Integrate with respect to \(y\) first.} Fix \(x\in(0,1)\). As
a function of \(y\in[0,1]\), the formula extends continuously and is
bounded, so this section is absolutely integrable. Direct
differentiation gives

\[
\frac{\partial}{\partial y}\left(\frac{y}{x^2+y^2}\right)
=\frac{x^2-y^2}{(x^2+y^2)^2}.
\]

Consequently

\[
\int_0^1 f(x,y)\,dy
=\left[\frac{y}{x^2+y^2}\right]_{y=0}^{y=1}
=\frac1{1+x^2},
\]

and the resulting outer function is integrable:

\[
\int_0^1\left(\int_0^1f(x,y)\,dy\right)dx
=\int_0^1\frac{dx}{1+x^2}
=\boxed{\frac\pi4}.
\]

\textbf{2. Integrate with respect to \(x\) first.} For each fixed
\(y\in(0,1)\), the section is again continuous and bounded on \([0,1]\).
This time

\[
\frac{\partial}{\partial x}\left(-\frac{x}{x^2+y^2}\right)
=\frac{x^2-y^2}{(x^2+y^2)^2}.
\]

Thus

\[
\int_0^1 f(x,y)\,dx=-\frac1{1+y^2},
\]

and

\[
\int_0^1\left(\int_0^1f(x,y)\,dx\right)dy
=-\int_0^1\frac{dy}{1+y^2}
=\boxed{-\frac\pi4}.
\]

Every section used here is integrable, not just almost every section.
That still does not establish joint absolute integrability.

\textbf{3. Prove both sign parts have infinite integral by explicit
disjoint rectangles.} If \(0<y<x/2\), then

\[
x^2-y^2\geq\frac34x^2,
\qquad (x^2+y^2)^2\leq\frac{25}{16}x^4,
\qquad f(x,y)\geq\frac{12}{25x^2}>0.
\]

For each integer \(n\geq0\), take

\[
R_n=\bigl(2^{-n-1},2^{-n}\bigr)
\times\bigl(0,2^{-n-2}\bigr).
\]

These rectangles are pairwise disjoint and lie in \(\{0<y<x/2\}\cap Q\).
On \(R_n\), \(x<2^{-n}\), so

\[
f^+(x,y)\geq\frac{12}{25}2^{2n},
\qquad \mathcal L^2(R_n)=2^{-2n-3}.
\]

For every positive integer \(N\), monotonicity of the nonnegative
integral yields

\[
\begin{aligned}
\int_Qf^+\,d\mathcal L^2
&\geq\sum_{n=0}^{N-1}
\frac{12}{25}2^{2n}\mathcal L^2(R_n)\\
&=\sum_{n=0}^{N-1}\frac3{50}
=\frac{3N}{50}.
\end{aligned}
\]

Letting \(N\to\infty\) proves \(\int_Qf^+=\infty\). Since
\(f(y,x)=-f(x,y)\), the transposed rectangles give the identical bound
for \(f^-\). Hence

\[
\int_Q f^+\,d\mathcal L^2
=\int_Q f^-\,d\mathcal L^2=\infty,
\qquad \int_Q|f|\,d\mathcal L^2=\infty.
\]

The expression \(\int_Qf\) is therefore \textbf{undefined as a signed
Lebesgue integral}, rather than equal to either iterated value or to
zero.

\textbf{4. Identify another prescription, without confusing it with a
Lebesgue integral.} On \((\varepsilon,1)^2\subset Q\), where
\(0<\varepsilon<1\), the function is bounded and antisymmetric under
exchanging coordinates. Fubini is valid on that truncated square, and

\[
\int_{(\varepsilon,1)^2}f\,d\mathcal L^2=0.
\]

Thus this particular symmetric limiting prescription gives zero. The
disagreement with the two iterated values records dependence on the
truncation procedure; it supplies no value for the undefined signed
Lebesgue integral.

\subsection{Why the theorem applies}\label{why-the-theorem-applies-3}

The one-dimensional computations use the fundamental theorem of calculus
on continuous sections. Fubini applies on each truncated square because
the integrand is bounded there. It does \textbf{not} apply on \(Q\):
joint absolute integrability fails by Step 3. Tonelli remains valid
separately for the nonnegative functions \(f^+\), \(f^-\), and \(|f|\),
with value \(\infty\) where indicated. It never legitimizes subtracting
\(\infty-\infty\).

\subsection{Pitfalls}\label{pitfalls-3}

\begin{itemize}
\tightlist
\item
  Finite iterated \textbf{signed} integrals do not replace the
  hypothesis \(\int_Q|f|<\infty\).
\item
  Cancellation inside each section can conceal infinite joint positive
  and negative masses.
\item
  Symmetry gives an integral of zero only after the existence of that
  integral has been justified, or for a specifically stated finite
  truncation.
\item
  Defining a value for \(f\) at the missing origin would not repair the
  divergence in its neighborhood.
\end{itemize}

\subsection{Reusable technique}\label{reusable-technique-3}

\begingroup\small\raggedright Tags: \texttt{signed-\allowbreak{}Fubini-\allowbreak{}failure}, \texttt{absolute-\allowbreak{}integrability},
\texttt{positive-\allowbreak{}negative-\allowbreak{}parts}, \texttt{dyadic-\allowbreak{}rectangles},
\texttt{cancellation}, \texttt{truncation-\allowbreak{}dependence}.\par\endgroup

Before exchanging signed integrals, estimate \(|f|\) or its sign parts.
A sequence of disjoint regions each carrying a fixed positive mass
proves divergence with an explicit sum, even when direct polar
integration is inconvenient. A sufficient repair is to impose
\(f\in L^1(Q)\); local boundedness away from the singularity is
insufficient.

\subsection{Sources}\label{sources-3}

\begin{itemize}
\tightlist
\item
  John K. Hunter,
  \href{https://www.math.ucdavis.edu/~hunter/measure_theory/measure_notes.pdf}{\emph{Measure
  Theory}}, Definition 4.8 (signed integral), Theorem 5.18 (Fubini), and
  Theorem 5.21 (completed Lebesgue measures). Source ID:
  \texttt{hunter-\allowbreak{}measure}. The primitive identities and dyadic-rectangle
  lower bounds are checked directly above.
\end{itemize}

\section{E05 --- Two finite iterated sums and no signed double
integral}\label{E05}

\subsection{Problem}\label{problem-4}

Let \(\mathbb N=\{1,2,\ldots\}\) and define the function
\(a:\mathbb N^2\to\mathbb R\) by \[
a(n,m)=\begin{cases}1&m=n,\\-1&m=n+1,\\0&\text{otherwise}.\end{cases}
\] Compute both iterated sums. Explain exactly why their disagreement
does not contradict Fubini's theorem, and what square partial sums do.

\subsection{Data and hypotheses}\label{data-and-hypotheses-4}

Each copy of \(\mathbb N\) has the power-set \(\sigma\)-algebra and
counting measure \(\#\); the product is counting measure on
\(\mathbb N^2\). These spaces are \(\sigma\)-finite because
\(\mathbb N\) is a countable union of finite sets. Every function here
is measurable. Write \(a_{nm}=a(n,m)\) for a \textbf{value} of the
function \(a\). A signed Lebesgue integral is defined only if at least
one of \(\int a^+\,d(\#\otimes\#)\) and \(\int a^-\,d(\#\otimes\#)\) is
finite; an integrable real-valued function must satisfy
\(\int|a|<\infty\).

\subsection{Solution}\label{solution-4}

\textbf{1. Sum each row first.} For every fixed \(n\), only two entries
are nonzero: \[
\sum_{m=1}^{\infty}a_{nm}=a_{nn}+a_{n,n+1}=1-1=0.
\] Consequently \[
\sum_{n=1}^{\infty}\left(\sum_{m=1}^{\infty}a_{nm}\right)=\sum_{n=1}^{\infty}0=0.
\] Every inner sum is absolutely convergent, and the resulting outer
series is absolutely convergent.

\textbf{2. Sum each column first.} Column \(m=1\) contains only
\(a_{11}=1\). For every \(m\ge2\), its nonzero entries are \(a_{mm}=1\)
and \(a_{m-1,m}=-1\). Thus \[
\sum_{n=1}^{\infty}a_{n1}=1,\qquad
\sum_{n=1}^{\infty}a_{nm}=0\quad(m\ge2),
\] and \[
\sum_{m=1}^{\infty}\left(\sum_{n=1}^{\infty}a_{nm}\right)=1+\sum_{m=2}^{\infty}0=1.
\] Again, every inner sum and the resulting outer series converge
absolutely. What fails is absolute summability \textbf{before}
cancellation: \[
\sum_{n=1}^{\infty}\sum_{m=1}^{\infty}|a_{nm}|
=\sum_{n=1}^{\infty}2=+\infty.
\] Indeed the positive diagonal and negative superdiagonal each contain
infinitely many entries of magnitude \(1\), so
\(\int a^+=\int a^-=+\infty\). The signed double integral
\(\int a\,d(\#\otimes\#)\) is undefined; neither iterated answer can be
identified with it.

\textbf{3. Inspect finite rectangles.} For positive integers \(N,M\), \[
S_{N,M}=\sum_{n=1}^{N}\sum_{m=1}^{M}a_{nm}
=\min(N,M)-\min(N,M-1).
\] Therefore \(S_{N,N}=1\), whereas \(S_{N,N+1}=0\). There is no limit
as \(N,M\to\infty\) independently. Taking \(M\to\infty\) first gives
\(0\); taking \(N\to\infty\) first gives \(1\). A square exhaustion
happens to select the second value, but it does not define a signed
Lebesgue integral.

\subsection{Why the theorem applies}\label{why-the-theorem-applies-4}

Tonelli's theorem applies to each of \(|a|\), \(a^+\) and \(a^-\),
because they are nonnegative measurable functions on \(\sigma\)-finite
product spaces. It correctly gives infinite integrals. Signed Fubini
requires \(\sum_{n,m}|a_{nm}|<\infty\), which has explicitly failed. If
instead a doubly indexed real or complex family \(b_{nm}\) satisfies
that bound, Fubini permits interchanging its sums and identifies both
with the product integral.

\subsection{Pitfalls}\label{pitfalls-4}

\begin{itemize}
\tightlist
\item
  Absolute convergence of every row, every column, and even both series
  of already-summed values does not establish
  \(\sum_{n,m}|a_{nm}|<\infty\).
\item
  This is an iterated-summation counterexample, not a conditionally
  Lebesgue-integrable function. For real or complex functions, Lebesgue
  integrability entails absolute integrability. Ordinary conditionally
  convergent series use an ordering, which counting-measure integration
  does not supply.
\item
  Never replace the undefined expression \(+\infty-(+\infty)\) by zero
  after pairing entries.
\end{itemize}

\subsection{Reusable technique}\label{reusable-technique-4}

Before exchanging signed sums or integrals, apply Tonelli to the
absolute value. If it diverges, inspect positive and negative parts and
the actual limiting convention. Tags: \texttt{counting-\allowbreak{}measure},
\texttt{signed-\allowbreak{}fubini}, \texttt{absolute-\allowbreak{}integrability},
\texttt{order-\allowbreak{}of-\allowbreak{}summation}, \texttt{counterexample}.

\subsection{Sources}\label{sources-4}

\begin{itemize}
\tightlist
\item
  John K. Hunter,
  \href{https://www.math.ucdavis.edu/~hunter/measure_theory/measure_notes.pdf}{\emph{Measure
  Theory}}, Theorem 5.18 and Example 5.19 (Fubini and absolutely
  convergent double series), and Chapter 4 (signed integration). Source
  ID: \texttt{hunter-\allowbreak{}measure}.
\item
  The matrix, explicit rectangular calculation, and exposition above
  were worked out for this collection; the source supports the theorem,
  not a copied exercise.
\end{itemize}

\section{E06 --- A logarithm from a positive double integral}\label{E06}

\subsection{Problem}\label{problem-5}

For \(a,b>0\), calculate \[
F(a,b)=\int_0^{\infty}\frac{e^{-ax}-e^{-bx}}{x}\,dx
\] as a Lebesgue integral. Justify both the exchange of integrations and
differentiation with respect to \(b\).

\subsection{Data and hypotheses}\label{data-and-hypotheses-5}

The integration variable \(x\) carries one-dimensional Lebesgue measure
\(\lambda^1\) on \((0,\infty)\). A parameter variable \(t\) carries
\(\lambda^1\) on a finite interval. The expression
\((e^{-ax}-e^{-bx})/x\) is the value at \(x\) of a continuous function
on \((0,\infty)\). Positivity of both parameters is essential for the
tail estimate below. No value at \(x=0\) is needed; assigning the
limiting value \(b-a\) there would not change the integral.

\subsection{Solution}\label{solution-5}

\textbf{1. Represent the cancellation before integrating.} Suppose first
\(a\le b\). The fundamental theorem of calculus in the parameter gives,
for \(x>0\), \[
\frac{e^{-ax}-e^{-bx}}{x}=\int_a^b e^{-tx}\,dt\ge0.
\] The function \((x,t)\mapsto e^{-tx}\) is continuous and nonnegative.
Tonelli therefore allows an exchange before finiteness has been
established: \[
\begin{aligned}
F(a,b)&=\int_0^{\infty}\int_a^b e^{-tx}\,dt\,dx\\
&=\int_a^b\int_0^{\infty}e^{-tx}\,dx\,dt
=\int_a^b\frac{1}{t}\,dt
=\log b-\log a.
\end{aligned}
\] The inner integral is \(1/t\) because \(t\ge a>0\). The resulting
logarithm is finite, so this also proves integrability. If \(a>b\),
apply this computation to \(F(b,a)\) and use the pointwise negative
sign; if \(a=b\), the integrand is identically zero. In all cases, \[
\boxed{F(a,b)=\log(b/a).}
\]

\textbf{2. Check absolute integrability directly.} Let \(c=\min(a,b)>0\)
and \(d=|b-a|\). The same parameter representation with the endpoints
ordered gives \[
\left|\frac{e^{-ax}-e^{-bx}}{x}\right|
\le d e^{-cx}\le d\qquad(x>0).
\] Thus the integral of the absolute value on \((0,1)\) is at most
\(d\). On \((1,\infty)\) one may also use \[
\left|\frac{e^{-ax}-e^{-bx}}{x}\right|
\le e^{-ax}+e^{-bx},
\] whose integral is finite. This makes the endpoint hypotheses visible
without relying on the final formula.

\textbf{3. Differentiate under the integral with a local bound.} Fix
\(a>0\) and \(b_0>0\), and write \(q_b(x)=(e^{-ax}-e^{-bx})/x\). Its
parameter derivative is \[
\partial_b q_b(x)=e^{-bx}.
\] For \(b\in[b_0/2,3b_0/2]\), \[
|\partial_b q_b(x)|\le e^{-b_0x/2},\qquad
\int_0^{\infty}e^{-b_0x/2}\,dx=2/b_0<\infty.
\] For sufficiently small nonzero \(h\), the ordinary mean value theorem
bounds the difference quotient \((q_{b_0+h}(x)-q_{b_0}(x))/h\) by this
same integrable function. These quotients converge pointwise to
\(e^{-b_0x}\). Dominated convergence yields \[
\partial_b F(a,b_0)=\int_0^{\infty}e^{-b_0x}\,dx=1/b_0,
\] in agreement with the logarithm. This is a proof of the exchange,
rather than a deduction from the already-known answer.

\subsection{Why the theorem applies}\label{why-the-theorem-applies-5}

Tonelli applies on the \(\sigma\)-finite Lebesgue product
\((0,\infty)\times[a,b]\) to a nonnegative Borel function. Once its
integral is finite, signed Fubini would also be available. For
differentiation, the difference quotients are measurable, converge
pointwise, and have a common \(L^1\) majorant on a neighborhood of the
chosen positive parameter; Hunter's dominated convergence theorem
applies. The same argument works near any \(b_0>0\); a uniform bound
over all \(b>0\) is unnecessary.

\subsection{Pitfalls}\label{pitfalls-5}

\begin{itemize}
\tightlist
\item
  The two terms \(e^{-ax}/x\) and \(e^{-bx}/x\) each have infinite
  integrals near zero. Writing their separate integrals and subtracting
  creates an undefined \(\infty-\infty\).
\item
  If \(a=0<b\), then \((1-e^{-bx})/x\) has a divergent tail, so the
  stated finite formula does not extend to \(a=0\).
\item
  A bound on \(q_b\) alone does not justify differentiating it. The
  proof controls its difference quotients through \(\partial_b q_b\).
\end{itemize}

\subsection{Reusable technique}\label{reusable-technique-5}

Turn a difference quotient or cancellation into an integral over a
positive parameter kernel. Apply Tonelli before evaluating, then use a
local integrable derivative bound for parameter differentiation. Tags:
\texttt{frullani}, \texttt{tonelli}, \texttt{parameter-\allowbreak{}integral},
\texttt{differentiate-\allowbreak{}under-\allowbreak{}integral}, \texttt{cancellation}.

\subsection{Sources}\label{sources-5}

\begin{itemize}
\tightlist
\item
  John K. Hunter,
  \href{https://www.math.ucdavis.edu/~hunter/measure_theory/measure_notes.pdf}{\emph{Measure
  Theory}}, Theorem 4.24 (dominated convergence) and Chapter 5 (product
  integration). Source ID: \texttt{hunter-\allowbreak{}measure}.
\item
  Keith Conrad,
  \href{https://kconrad.math.uconn.edu/blurbs/analysis/diffunderint.pdf}{\emph{Differentiating
  under the integral sign}}, Theorem 12.3, for a classical sufficient
  differentiation theorem. The measurable version used here follows from
  the displayed difference-quotient proof. Source ID:
  \texttt{conrad-\allowbreak{}diff}.
\end{itemize}

\section{E07 --- A weighted integral on a parabolic arc}\label{E07}

\subsection{Problem}\label{problem-6}

Let \(C=\{(t,t^2):0\leq t\leq1\}\). Compute the integral of the
coordinate function \(f:C\to\mathbb R\), defined by \(f(x,y)=x\), with
respect to arc length. Also compute the length of \(C\).

\subsection{Data and hypotheses}\label{data-and-hypotheses-6}

The parametrization is \(\gamma:[0,1]\to\mathbb R^2\),
\(\gamma(t)=(t,t^2)\). The curve \(C\) is compact, and \(f\) is
continuous and nonnegative on \(C\).

Throughout, \(\mathcal L^k\) denotes Lebesgue measure. For an integer
\(k\geq1\), use the normalized Hausdorff measure

\[
\mathcal H^k(A)=\lim_{\delta\downarrow0}\inf\left\{\sum_{j=1}^{\infty}\frac{\alpha_k}{2^k}(\operatorname{diam}U_j)^k:
A\subseteq\bigcup_{j=1}^{\infty}U_j,\ \operatorname{diam}U_j<\delta\right\},
\]

where \(\alpha_k=\mathcal L^k(B^k(0,1))\). Thus \(\mathcal H^k\) on a
\(k\)-plane agrees with \(k\)-dimensional Lebesgue measure, and
\(\mathcal H^0\) is counting measure. In particular, \(\alpha_1=2\), so
\(\mathcal H^1\) here is ordinary arc length. The measure being
integrated is the restriction \(\mathcal H^1\mathbin{\vert}C\), not
two-dimensional area measure.

\subsection{Solution}\label{solution-6}

Differentiate the parametrization:

\[
\gamma'(t)=(1,2t),\qquad J_1\gamma(t)=|\gamma'(t)|=\sqrt{1+4t^2}.
\]

The first coordinate proves injectivity: if \(\gamma(s)=\gamma(t)\),
then \(s=t\). The weighted area formula therefore gives

\[
\int_C f\,d\mathcal H^1
=\int_0^1 f(\gamma(t))|\gamma'(t)|\,dt
=\int_0^1 t\sqrt{1+4t^2}\,dt.
\]

Here \(f\) is the function; \(f(\gamma(t))=t\) is its value at the point
\(\gamma(t)\). With \(u=1+4t^2\) and \(du=8t\,dt\),

\[
\int_0^1 t\sqrt{1+4t^2}\,dt
=\frac18\int_1^5 u^{1/2}\,du
=\frac1{12}[u^{3/2}]_1^5
=\boxed{\frac{5\sqrt5-1}{12}}.
\]

For length, replace \(f\) by the constant function \(1\):

\[
\mathcal H^1(C)=\int_0^1\sqrt{1+4t^2}\,dt.
\]

Put \(s=2t\). The identity

\[
\frac{d}{ds}\left(\frac12\bigl(s\sqrt{1+s^2}+\operatorname{arsinh}s\bigr)\right)
=\sqrt{1+s^2}
\]

yields

\[
\mathcal H^1(C)
=\frac{\sqrt5}{2}+\frac14\operatorname{arsinh}2
=\boxed{\frac{\sqrt5}{2}+\frac14\log(2+\sqrt5)}.
\]

\subsection{Why the theorem applies}\label{why-the-theorem-applies-6}

The map \(\gamma\) is \(C^1\) on a neighborhood of \([0,1]\) and has
bounded derivative there, hence is Lipschitz on this interval. Its
domain is Borel. Its multiplicity is exactly one on \(C\). The weighted
area formula applies to the nonnegative Borel weight \(f\circ\gamma\);
the explicit integral is finite since
\(0\leq t\sqrt{1+4t^2}\leq\sqrt5\). Taking the constant weight also
proves that \(C\) has finite \(\mathcal H^1\) measure.

\subsection{Pitfalls}\label{pitfalls-6}

\begin{itemize}
\tightlist
\item
  Integrating \(t\,dt\) gives \(1/2\), but this uses horizontal distance
  instead of arc length. The speed factor is necessary.
\item
  The expression \(\int_C f\,d\mathcal L^2\) would be zero: \(C\) has
  two-dimensional Lebesgue measure zero. A Lipschitz arc can be covered
  by \(N\) disks of radius proportional to \(1/N\), with total area
  tending to zero.
\item
  Reversing the traversal leaves this scalar Hausdorff integral
  unchanged. A vector circulation would include a signed tangent vector
  and can change sign.
\end{itemize}

\subsection{Reusable technique}\label{reusable-technique-6}

For an injective Lipschitz curve, compose the scalar weight with the
parametrization, then multiply by speed. Coordinate weights often turn a
square-root integral into a direct substitution.

\begingroup\small\raggedright Tags: \texttt{hausdorff-\allowbreak{}normalization}, \texttt{weighted-\allowbreak{}arc-\allowbreak{}length},
\texttt{injective-\allowbreak{}area-\allowbreak{}formula}, \texttt{substitution},
\texttt{coordinate-\allowbreak{}moment}.\par\endgroup

\subsection{Sources}\label{sources-6}

\begin{itemize}
\tightlist
\item
  Leon Simon,
  \href{https://math.stanford.edu/~lms/ntu-gmt-text.pdf}{\emph{Introduction
  to Geometric Measure Theory} (2018)}, Chapter 1, §2, definitions
  2.1--2.2 and Theorem 2.6 (normalization); Chapter 2, §3, Theorem 3.3,
  Remark 3.4(2), and Example 3.5(1) (weighted area formula and arc
  length). The specific curve and computations above are worked out
  here. Source ID: \texttt{simon-\allowbreak{}gmt}.
\end{itemize}

\section{E08 --- Traversing a circle twice: multiplicity and
orientation}\label{E08}

\subsection{Problem}\label{problem-7}

Let \(C=\{(x,y)\in\mathbb R^2:x^2+y^2=1\}\) and let \(w:C\to\mathbb R\)
be the function \(w(x,y)=x^2\). Compare

\[
\int_C w\,d\mathcal H^1
\quad\hbox{and}\quad
\int_0^{4\pi}w(\gamma(t))|\gamma'(t)|\,dt,
\qquad \gamma(t)=(\cos t,\sin t).
\]

Then compare these scalar integrals with the circulation of the vector
field \(F(x,y)=(-y,x)\).

\subsection{Data and hypotheses}\label{data-and-hypotheses-7}

Use normalized Hausdorff measure: \(\mathcal H^k\) has covering cost
\(\alpha_k(\operatorname{diam}U/2)^k\), where \(\alpha_k\) is the volume
of the unit ball in \(\mathbb R^k\). It agrees with \(\mathcal L^k\) on
\(k\)-planes; \(\mathcal H^0\) counts points. Thus
\(\mathcal H^1\mathbin{\vert}C\) is ordinary length measure. The
functions \(w\) and \(F\) are continuous, and \(0\leq w\leq1\) on \(C\).

The map \(\gamma\) is Lipschitz on \(A=[0,4\pi]\). Its image is the
geometric set \(C\), which has no record of how many times it was
traversed.

\subsection{Solution}\label{solution-7}

First compute the parameter integral. Since

\[
\gamma'(t)=(-\sin t,\cos t),\qquad |\gamma'(t)|=1,
\qquad w(\gamma(t))=\cos^2t,
\]

we obtain

\[
\int_0^{4\pi}w(\gamma(t))|\gamma'(t)|\,dt
=\int_0^{4\pi}\frac{1+\cos(2t)}2\,dt
=\boxed{2\pi}.
\]

This is a traversal integral. To understand its relation to the
geometric integral, define

\[
N(q)=\mathcal H^0\bigl(\{t\in A:\gamma(t)=q\}\bigr).
\]

If \(q\in C\setminus\{(1,0)\}\), there are exactly two preimages, one in
\((0,2\pi)\) and one in \((2\pi,4\pi)\). At \((1,0)\) the preimages are
\(0,2\pi,4\pi\), so \(N(1,0)=3\). Off \(C\), \(N=0\). The single
exceptional image point has \(\mathcal H^1\) measure zero. The weighted
area formula therefore reads

\[
\begin{aligned}
\int_0^{4\pi}w(\gamma(t))|\gamma'(t)|\,dt
&=\int_C\sum_{t\in\gamma^{-1}(q)\cap A}w(q)\,d\mathcal H^1(q)\\
&=\int_C w(q)N(q)\,d\mathcal H^1(q)
=2\int_C w\,d\mathcal H^1.
\end{aligned}
\]

Consequently,

\[
\boxed{\int_C w\,d\mathcal H^1=\pi}.
\]

A single traversal \(0\leq t\leq2\pi\) gives the same geometric answer
directly: only the common endpoint is counted twice, and it is
length-null.

For circulation, the integrand includes the velocity vector rather than
its norm:

\[
F(\gamma(t))=(-\sin t,\cos t)=\gamma'(t),
\qquad F(\gamma(t))\cdot\gamma'(t)=1.
\]

Thus the double counterclockwise traversal has circulation

\[
\int_0^{4\pi}F(\gamma(t))\cdot\gamma'(t)\,dt=\boxed{4\pi},
\]

and a single counterclockwise traversal has circulation \(2\pi\). A
reversed single traversal has circulation \(-2\pi\), whereas its scalar
\(w\)-integral remains \(\pi\). These circulation values are not
integrals of the scalar function \(w\).

Flux requires a chosen unit normal: it has the form
\(\int_C G\cdot n\,d\mathcal H^1\). For example, \(G(x,y)=(x,y)\) has
outward flux \(2\pi\), since the outward normal is \(n(q)=q\) and
\(G(q)\cdot n(q)=1\). Choosing the inward normal changes that flux to
\(-2\pi\). The positive measure \(\mathcal H^1\) itself supplies neither
a tangent orientation nor a normal orientation.

\subsection{Why the theorem applies}\label{why-the-theorem-applies-7}

The domain \(A\) is compact Borel and \(\gamma\) is \(C^1\) with bounded
derivative. The nonnegative weight \(t\mapsto w(\gamma(t))\) is Borel.
Hence the weighted area formula applies even though \(\gamma\) is not
injective. Its fiber sum is explicit above. Integrability is also
direct: the parameter integrand is bounded by \(1\) on a finite
interval. The same formula with weight \(1\) shows
\(\mathcal H^1(C)=2\pi\), which justifies all bounded scalar and flux
integrals used here.

\subsection{Pitfalls}\label{pitfalls-7}

\begin{itemize}
\tightlist
\item
  A parametrization's image alone does not determine its traversal
  integral. Multiplicity must be included.
\item
  Multiplicity is a count, never a signed count in the scalar area
  formula. Opposite traversals add for scalar length even when their
  circulations cancel.
\item
  Endpoint multiplicity is genuinely three at \((1,0)\) for the stated
  closed domain. Ignoring it is justified by its null image, not by
  claiming it is two everywhere.
\end{itemize}

\subsection{Reusable technique}\label{reusable-technique-7}

Before using a noninjective parametrization, count its fibers. The area
formula converts a parameter integral into an integral of an explicit
sum over preimages. Treat orientation separately whenever a vector field
is present.

\begingroup\small\raggedright Tags: \texttt{area-\allowbreak{}formula}, \texttt{multiplicity},
\texttt{noninjective-\allowbreak{}parametrization}, \texttt{scalar-\allowbreak{}versus-\allowbreak{}vector},
\texttt{circulation}, \texttt{flux}.\par\endgroup

\subsection{Sources}\label{sources-7}

\begin{itemize}
\tightlist
\item
  Leon Simon,
  \href{https://math.stanford.edu/~lms/ntu-gmt-text.pdf}{\emph{Introduction
  to Geometric Measure Theory} (2018)}, Chapter 1, §2, definitions
  2.1--2.2 and Theorem 2.6; Chapter 2, §3, Theorem 3.3 and Remark 3.4(2)
  (weighted fiber sum). The fiber count and vector comparisons above are
  direct computations. Source ID: \texttt{simon-\allowbreak{}gmt}.
\end{itemize}

\section{E09 --- Height integral over a paraboloid graph}\label{E09}

\subsection{Problem}\label{problem-8}

Let \(D=\{(x,y)\in\mathbb R^2:x^2+y^2\leq1\}\) and

\[
S=\{(x,y,x^2+y^2):(x,y)\in D\}.
\]

For the height function \(h:S\to\mathbb R\), \(h(x,y,z)=z\), compute
\(\int_S h\,d\mathcal H^2\). Also compute the area of \(S\).

\subsection{Data and hypotheses}\label{data-and-hypotheses-8}

All measures are positive scalar measures. The normalized Hausdorff
measure has covering cost \(\alpha_k(\operatorname{diam}U/2)^k\), with
\(\alpha_k=\mathcal L^k(B^k(0,1))\); it agrees with \(\mathcal L^k\) on
\(k\)-planes. Thus the cost for \(\mathcal H^2\) is
\((\pi/4)(\operatorname{diam}U)^2\), and \(\mathcal H^0\) is counting
measure.

Write \(u:D\to\mathbb R\), \(u(x,y)=x^2+y^2\), and
\(\Phi:D\to\mathbb R^3\), \(\Phi(x,y)=(x,y,u(x,y))\). The domain \(D\)
is compact Borel. The map \(\Phi\) is injective because its first two
coordinates recover its input. The function \(h\) is continuous with
\(0\leq h\leq1\) on \(S\).

\subsection{Solution}\label{solution-8}

Compute the derivative columns and their Gram matrix:

\[
\partial_x\Phi=(1,0,2x),\qquad
\partial_y\Phi=(0,1,2y),
\]

\[
(D\Phi)^T D\Phi=
\begin{pmatrix}1+4x^2&4xy\\4xy&1+4y^2\end{pmatrix}.
\]

Its determinant is

\[
(1+4x^2)(1+4y^2)-16x^2y^2=1+4x^2+4y^2.
\]

Hence the two-dimensional area factor is

\[
J_2\Phi(x,y)=\sqrt{1+4x^2+4y^2}
=\sqrt{1+|\nabla u(x,y)|^2}.
\]

Injectivity and the weighted area formula give

\[
\int_S h\,d\mathcal H^2
=\int_D h(\Phi(x,y))J_2\Phi(x,y)\,d\mathcal L^2(x,y)
=\int_D(x^2+y^2)\sqrt{1+4x^2+4y^2}\,dx\,dy.
\]

The function \(h\) is defined on \(S\); its pulled-back value is
\(h(\Phi(x,y))=u(x,y)\). Now use polar coordinates in the planar disk,
including their separate Jacobian factor \(r\):

\[
I=2\pi\int_0^1 r^3\sqrt{1+4r^2}\,dr.
\]

For \(v=1+4r^2\), we have \(r\,dr=dv/8\) and \(r^2=(v-1)/4\), so

\[
\begin{aligned}
I&=\frac{\pi}{16}\int_1^5(v-1)v^{1/2}\,dv\\
&=\frac{\pi}{16}\left[\frac25v^{5/2}-\frac23v^{3/2}\right]_1^5\\
&=\frac{\pi}{16}\left(\frac{20\sqrt5}{3}+\frac4{15}\right)
=\boxed{\frac{\pi(25\sqrt5+1)}{60}}.
\end{aligned}
\]

Using weight \(1\) instead gives

\[
\begin{aligned}
\mathcal H^2(S)
&=2\pi\int_0^1r\sqrt{1+4r^2}\,dr\\
&=\frac\pi4\int_1^5v^{1/2}\,dv
=\boxed{\frac{\pi(5\sqrt5-1)}6}.
\end{aligned}
\]

As a quick consistency check, \(0<I<\mathcal H^2(S)\), as expected
because the height lies between \(0\) and \(1\) and is strictly between
them on a set of positive area.

\subsection{Why the theorem applies}\label{why-the-theorem-applies-8}

On the convex disk, \(|\nabla u|\leq2\), so \(|u(p)-u(q)|\leq2|p-q|\)
and \(|\Phi(p)-\Phi(q)|\leq\sqrt5|p-q|\). Thus \(\Phi\) is Lipschitz on
\(D\). Its image is compact Borel, its multiplicity is one, and the
weight is nonnegative Borel. The explicit area factor is bounded by
\(\sqrt5\), so both pulled-back integrands are integrable on \(D\). The
polar coordinate change is applied to a bounded, nonnegative measurable
function; its seam and origin have planar measure zero. The area formula
handles the closed disk directly, so there is no need to assume an open
parameter domain or discard the boundary without explanation.

\subsection{Pitfalls}\label{pitfalls-8}

\begin{itemize}
\tightlist
\item
  The projected integral \(\int_D u\,dx\,dy=\pi/2\) differs from the
  surface integral. Projection omits the graph's area factor.
\item
  There are two distinct factors after the two changes of variables:
  \(\sqrt{1+4r^2}\) comes from graph area; \(r\) comes from planar polar
  coordinates.
\item
  The cross product \((-2x,-2y,1)\) can also compute the area factor,
  but only its norm enters this scalar integral. A signed vector flux
  would additionally require a chosen normal.
\end{itemize}

\subsection{Reusable technique}\label{reusable-technique-8}

For a graph \(z=u(x,y)\) with one point above each base point, pull the
scalar weight back and multiply by \(\sqrt{1+|\nabla u|^2}\). If the
domain and pulled-back integrand are radial, use polar coordinates only
after forming this graph factor.

\begingroup\small\raggedright Tags: \texttt{graph-\allowbreak{}area}, \texttt{weighted-\allowbreak{}surface-\allowbreak{}integral},
\texttt{gram-\allowbreak{}determinant}, \texttt{polar-\allowbreak{}coordinates},
\texttt{injective-\allowbreak{}area-\allowbreak{}formula}, \texttt{substitution}.\par\endgroup

\subsection{Sources}\label{sources-8}

\begin{itemize}
\tightlist
\item
  Leon Simon,
  \href{https://math.stanford.edu/~lms/ntu-gmt-text.pdf}{\emph{Introduction
  to Geometric Measure Theory} (2018)}, Chapter 1, §2, definitions
  2.1--2.2 and Theorem 2.6; Chapter 2, §3, definition 3.2, Theorem 3.3,
  Remark 3.4(2), and Example 3.5(3) (graph area factor). The Gram
  determinant, substitutions, and numerical constants above are derived
  here. Source ID: \texttt{simon-\allowbreak{}gmt}.
\end{itemize}

\section{E10 --- A sphere's second moment and a spherical
cap}\label{E10}

\subsection{Problem}\label{problem-9}

For \(R>0\), let \(S_R=\{(x,y,z)\in\mathbb R^3:x^2+y^2+z^2=R^2\}\).
Compute the scalar integral of the function \(h:S_R\to\mathbb R\),
\(h(x,y,z)=z^2\):

\[
\int_{S_R}h\,d\mathcal H^2.
\]

As a related application, compute the area of a cap of height
\(a\in[0,2R]\) defined by \(z\geq R-a\).

\subsection{Data and hypotheses}\label{data-and-hypotheses-9}

Use normalized Hausdorff measure: the covering cost is
\(\alpha_k(\operatorname{diam}U/2)^k\), where \(\alpha_k\) is the
\(k\)-dimensional unit-ball volume. This makes \(\mathcal H^k\) equal to
\(\mathcal L^k\) on \(k\)-planes; \(\mathcal H^0\) is counting measure.
Thus \(\mathcal H^2\) has covering cost
\((\pi/4)(\operatorname{diam}U)^2\) and is usual surface area on a
smooth surface. The sphere and cap are compact Borel sets, and \(h\) is
continuous and nonnegative.

\subsection{Solution}\label{solution-9}

Use the spherical map

\[
\Phi(\theta,\varphi)
=(R\sin\theta\cos\varphi,R\sin\theta\sin\varphi,R\cos\theta),
\]

on \(D=(0,\pi)\times(0,2\pi)\). Its derivative columns are
perpendicular, with lengths \(R\) and \(R\sin\theta\). Therefore

\[
J_2\Phi(\theta,\varphi)=R^2\sin\theta.
\]

The map is injective on \(D\): \(z/R\) determines \(\theta\in(0,\pi)\),
and the first two coordinates then determine the unique
\(\varphi\in(0,2\pi)\). Its image omits exactly the meridian

\[
M=\{(R\sin\theta,0,R\cos\theta):0\leq\theta\leq\pi\},
\]

including the poles. This omitted set is \(\mathcal H^2\)-null. Indeed,
its displayed curve is \(R\)-Lipschitz; partitioning \([0,\pi]\) into
\(N\) intervals covers \(M\) by \(N\) sets of diameter at most
\(R\pi/N\). Their total normalized two-dimensional cost is at most
\(\pi^3R^2/(4N)\), tending to zero, while their maximum diameter also
tends to zero.

The weighted area formula now gives

\[
\begin{aligned}
\int_{S_R}h\,d\mathcal H^2
&=\int_0^{2\pi}\int_0^\pi
R^2\cos^2\theta\,R^2\sin\theta\,d\theta\,d\varphi\\
&=2\pi R^4\int_0^\pi\cos^2\theta\sin\theta\,d\theta.
\end{aligned}
\]

With \(v=\cos\theta\), \(dv=-\sin\theta\,d\theta\),

\[
\int_0^\pi\cos^2\theta\sin\theta\,d\theta
=\int_{-1}^1v^2\,dv=\frac23,
\]

so

\[
\boxed{\int_{S_R}z^2\,d\mathcal H^2=\frac{4\pi R^4}{3}}.
\]

The same parametrization with weight \(1\) gives
\(\mathcal H^2(S_R)=4\pi R^2\). A separate check uses coordinate
symmetry: Hausdorff measure is unchanged by rotations, so the integrals
of \(x^2\), \(y^2\), and \(z^2\) are equal. Their sum is

\[
\int_{S_R}(x^2+y^2+z^2)\,d\mathcal H^2
=R^2\mathcal H^2(S_R)=4\pi R^4.
\]

For the cap, set \(\theta_0=\arccos(1-a/R)\). The condition
\(z\geq R-a\) is \(0\leq\theta\leq\theta_0\). Thus

\[
\mathcal H^2(\{q\in S_R:z(q)\geq R-a\})
=2\pi R^2\int_0^{\theta_0}\sin\theta\,d\theta
=\boxed{2\pi Ra}.
\]

This includes \(a=0\) (a single point, area zero) and \(a=2R\) (the
entire sphere).

\subsection{Why the theorem applies}\label{why-the-theorem-applies-9}

The map \(\Phi\) extends smoothly to the closed rectangle
\([0,\pi]\times[0,2\pi]\), with bounded derivative, hence is Lipschitz
there and on \(D\). The restricted domain \(D\) is Borel and has
multiplicity one on its image. Its Jacobian is bounded by \(R^2\), while
\(h\circ\Phi\) is bounded by \(R^2\) on a finite-measure rectangle. All
scalar integrals are therefore finite. The missing meridian is null by
the covering argument above. Nonnegative integrands justify the iterated
integrations by Tonelli; finiteness also permits Fubini. Cap boundary
latitudes are Lipschitz curves and hence \(\mathcal H^2\)-null by the
same covering argument.

\subsection{Pitfalls}\label{pitfalls-9}

\begin{itemize}
\tightlist
\item
  The factor \(R^2\sin\theta\) is essential. The angular parameter
  rectangle carries planar measure, not surface measure.
\item
  The closed-rectangle parametrization is not injective at the seam or
  poles. Restricting to \(D\) and proving that the omitted image is null
  resolves this explicitly.
\item
  This scalar moment has no chosen orientation. For a vector field
  \(G(q)=q\), the outward unit normal is \(n(q)=q/R\) and the outward
  flux is
  \(\int_{S_R}G\cdot n\,d\mathcal H^2=R\mathcal H^2(S_R)=4\pi R^3\). The
  inward normal reverses that flux, while the \(z^2\) integral remains
  unchanged.
\end{itemize}

\subsection{Reusable technique}\label{reusable-technique-9}

Compute a surface Jacobian from the Gram determinant or perpendicular
tangent vectors. Separate chart overlap/null-image issues from the
calculation. Use rotation symmetry to check coordinate moments, and
express caps through their polar angle or height.

\begingroup\small\raggedright Tags: \texttt{sphere}, \texttt{surface-\allowbreak{}moment}, \texttt{spherical-\allowbreak{}cap},
\texttt{area-\allowbreak{}formula}, \texttt{null-\allowbreak{}seam}, \texttt{rotation-\allowbreak{}symmetry},
\texttt{scalar-\allowbreak{}versus-\allowbreak{}flux}.\par\endgroup

\subsection{Sources}\label{sources-9}

\begin{itemize}
\tightlist
\item
  Leon Simon,
  \href{https://math.stanford.edu/~lms/ntu-gmt-text.pdf}{\emph{Introduction
  to Geometric Measure Theory} (2018)}, Chapter 1, §2, definitions
  2.1--2.2 and Theorem 2.6; Chapter 2, §3, definition 3.2, Theorem 3.3,
  Remark 3.4(2), and Example 3.5(2) (surface measure and area formula).
  The chart, null-meridian proof, moment, cap computation, and symmetry
  check above are derived here. Source ID: \texttt{simon-\allowbreak{}gmt}.
\end{itemize}

\section{E11 --- A fold counts the image twice}\label{E11}

\subsection{Problem}\label{problem-10}

Let \(A=[-1,1]\times[0,1]\) and define the map \(F:A\to\mathbb R^2\) by
\[
F(u,v)=(u^2,v).
\] Compute the image area, the integral of the area Jacobian over \(A\),
and the weighted integral \[
\int_A (u^2+v)J_2F(u,v)\,d\mathcal L^2(u,v).
\] Explain why an injective change-of-variables formula would give the
wrong answer. Then recover the scalar integral of \(h(x,y)=x+y\) over
the image, counted once.

\subsection{Data and hypotheses}\label{data-and-hypotheses-10}

The domain and target measures are completed two-dimensional Lebesgue
measure \(\mathcal L^2\). We use normalized Hausdorff measure \[
\mathcal H^k(E)=\lim_{\delta\downarrow0}\inf\left\{
\sum_j\frac{\alpha_k}{2^k}(\operatorname{diam}U_j)^k:
E\subseteq\bigcup_jU_j,\ \operatorname{diam}U_j<\delta\right\},
\] where \(\alpha_k\) is the volume of the unit ball in \(\mathbb R^k\).
Thus \(\mathcal H^k\) agrees with \(\mathcal L^k\) on a \(k\)-plane;
\(\mathcal H^0\) is counting measure.

The set \(A\) is compact and Borel. The map \(F\) is the restriction of
a smooth map, and is Lipschitz on \(A\): \(|u^2-s^2|\leq2|u-s|\) for
\(u,s\in[-1,1]\) gives \(|F(u,v)-F(s,t)|\leq2|(u,v)-(s,t)|\). The scalar
function \(h\) is continuous and nonnegative on the image. All integrals
computed below are finite. The multiplicity is the function \[
N(F,A,z)=\mathcal H^0\bigl(A\cap F^{-1}(\{z\})\bigr),
\] not the map \(F\) itself and not a chosen inverse branch.

\subsection{Solution}\label{solution-10}

\textbf{1. Compute the image and every fibre.} The image is
\(Q=[0,1]\times[0,1]\). For \(z=(x,y)\in Q\) with \(x>0\), the fibre
consists of the two distinct points \[
A\cap F^{-1}(\{(x,y)\})=\{(\sqrt x,y),(-\sqrt x,y)\}.
\] At \(x=0\) it consists of the single point \((0,y)\), and outside
\(Q\) it is empty. Consequently \(N=2\) on \((0,1]\times[0,1]\), \(N=1\)
on \(\{0\}\times[0,1]\), and \(N=0\) elsewhere. In particular \(N=2\) at
\(\mathcal L^2\)-almost every point of \(Q\).

\textbf{2. Compute the area Jacobian.} Differentiating the map gives \[
DF(u,v)=\begin{pmatrix}2u&0\\0&1\end{pmatrix},\qquad
J_2F(u,v)=\sqrt{\det(DF^TDF)}=|\det DF|=2|u|.
\] The critical set is \(\{0\}\times[0,1]\), where the rank drops to
one. Its image is the left edge of \(Q\); both sets have
\(\mathcal L^2\)-measure zero. We do not need to define an inverse
derivative there.

\textbf{3. Compare image area with the parameter integral.} The image
area is \[
\mathcal H^2(F(A))=\mathcal L^2(Q)=1.
\] Direct integration, justified by Tonelli for the nonnegative Borel
integrand, gives \[
\int_A J_2F\,d\mathcal L^2
=\int_{-1}^1\int_0^1 2|u|\,dv\,du
=4\int_0^1u\,du=2.
\] The area formula explains the factor: \[
\int_A J_2F\,d\mathcal L^2
=\int_{\mathbb R^2}N(F,A,z)\,d\mathcal H^2(z)=2\mathcal L^2(Q).
\] The extra unit is an entire second sheet, not an error caused by a
negligible boundary.

\textbf{4. Keep the fibre sum in the weighted formula.} Write the source
weight as \(h\circ F\), whose value at \((u,v)\) is \(u^2+v\). The
weighted area formula gives \[
\begin{aligned}
\int_A(h\circ F)J_2F\,d\mathcal L^2
&=\int_{\mathbb R^2}
\sum_{(u,v)\in A\cap F^{-1}(\{z\})}h(F(u,v))\,d\mathcal H^2(z)\\
&=\int_Qh(z)N(F,A,z)\,d\mathcal L^2(z)\\
&=2\int_0^1\int_0^1(x+y)\,dy\,dx=2.
\end{aligned}
\] An independent direct calculation is \[
\begin{aligned}
\int_{-1}^1\int_0^1(u^2+v)2|u|\,dv\,du
&=4\int_0^1\left(u^3+\frac u2\right)du\\
&=4\left(\frac14+\frac14\right)=2.
\end{aligned}
\] Thus a substitution that simply writes this parameter integral as
\(\int_Qh\) loses multiplicity.

\textbf{5. Count the image once.} Restrict \(F\) to
\(A_+=[0,1]\times[0,1]\). This restriction is injective, including its
boundary. Its Jacobian may vanish at \(u=0\); injectivity and Lipschitz
regularity are still enough for the area formula. It follows that \[
\int_Qh\,d\mathcal H^2
=\int_0^1\int_0^1(u^2+v)2u\,dv\,du=1.
\] Alternatively, divide the original parameter integral by two:
\(N(F,A,F(u,v))=2\) away from \(u=0\), and that exceptional line is
null. For this particular map, \[
\int_Qh\,d\mathcal H^2=\frac12\int_A(h\circ F)J_2F\,d\mathcal L^2=1.
\] A constant divisor works here because the multiplicity is constant
almost everywhere on the image. Variable multiplicity requires the
actual fibre count.

\subsection{Why the theorem applies}\label{why-the-theorem-applies-10}

The area formula applies to a Lipschitz map on a Lebesgue-measurable
domain, with nonnegative measurable weights. Its target integral
contains the sum over all preimages in that domain. Those hypotheses
hold here, so absolute integrability need not be assumed before
obtaining the nonnegative identity. Finiteness then follows from the
explicit computations. For a signed weight one must separately check
integrability, for example \(\int_A|h\circ F|J_2F<\infty\), before
subtracting positive and negative parts.

\subsection{Pitfalls}\label{pitfalls-10}

\begin{itemize}
\tightlist
\item
  The unsigned Jacobian is \(2|u|\), not \(2u\). Indeed
  \(\int_A\det DF\,d\mathcal L^2=0\) by oddness: opposite orientations
  cancel. Scalar Hausdorff area is positive and does not cancel.
\item
  Scalar Hausdorff integration supplies no normal vector. Vector flux
  requires an oriented normal and a dot product, in addition to the
  surface measure.
\item
  Noninjectivity cannot be dismissed by deleting only the critical line:
  the two regular sheets still have the same image.
\item
  The inverse branch \(x\mapsto\sqrt x\) has an unbounded derivative
  near zero. That does not invalidate this forward Lipschitz area
  calculation.
\end{itemize}

\subsection{Reusable technique}\label{reusable-technique-10}

Find the fibres before changing variables. Compute \(J_kF\) using
\(\sqrt{\det(DF^TDF)}\); keep the explicit preimage sum; then either
choose an injective sheet or correct for verified multiplicity. Tags:
\texttt{area-\allowbreak{}formula}, \texttt{multiplicity}, \texttt{fold},
\texttt{noninjective-\allowbreak{}parametrization}, \texttt{unsigned-\allowbreak{}jacobian},
\texttt{orientation}.

\subsection{Sources}\label{sources-10}

\begin{itemize}
\tightlist
\item
  Leon Simon,
  \href{https://math.stanford.edu/~lms/ntu-gmt-text.pdf}{\emph{Introduction
  to Geometric Measure Theory}}, Chapter 2, §3, Theorem 3.3 and Remark
  3.4(2), printed pp.~57--58; source ID \texttt{simon-\allowbreak{}gmt}. The
  same-dimensional specialization used here has target measure
  \(\mathcal H^2=\mathcal L^2\). The worked fold and calculations above
  are original to this collection.
\end{itemize}

\section{E12 --- Coarea for squared distance, including the critical
value}\label{E12}

\subsection{Problem}\label{problem-11}

Let \(D=\{(x,y)\in\mathbb R^2:x^2+y^2\leq1\}\) and define
\(g:\mathbb R^2\to\mathbb R\) by \(g(x,y)=x^2+y^2\). For a nonnegative
Borel function \(\phi:[0,1]\to[0,\infty]\), prove \[
\int_D\phi(g(x,y))\,d\mathcal L^2(x,y)
=\pi\int_0^1\phi(t)\,dt.
\] Use coarea, making the treatment of \(g^{-1}(\{0\})\) explicit. Also
compute \(\int_D|\nabla g|\,d\mathcal L^2\) and determine exactly when
the radial singularity \((x^2+y^2)^a\) is integrable.

\subsection{Data and hypotheses}\label{data-and-hypotheses-11}

Planar integration uses completed Lebesgue measure \(\mathcal L^2\). The
level sets carry normalized \(\mathcal H^1\), so their measure is
ordinary arc length. In general our normalization uses the covering cost
\(\alpha_k(\operatorname{diam}U)^k/2^k\), with \(\alpha_k\) the unit
\(k\)-ball volume; it gives \(\mathcal H^k=\mathcal L^k\) on a
\(k\)-plane and \(\mathcal H^0\) equal to counting measure.

The map \(g\) is smooth on \(\mathbb R^2\), the disk is Borel, and its
one-dimensional coarea Jacobian is \[
J_1g(x,y)=|\nabla g(x,y)|=2\sqrt{x^2+y^2}.
\] The weight \(\phi\) may be unbounded, and either side of the
requested formula may be \(+\infty\). No signed cancellation is
permitted in this extended nonnegative statement. A value assigned to a
function at the single point \((0,0)\) will not change its
\(\mathcal L^2\) integral.

\subsection{Solution}\label{solution-11}

\textbf{1. Describe the levels and the critical set.} Put
\(r=\sqrt{x^2+y^2}\). For \(0<t\leq1\) the level in \(D\) is the circle
of radius \(\sqrt t\): \[
D\cap g^{-1}(\{t\})=\{(\sqrt t\cos\theta,\sqrt t\sin\theta):0\leq\theta<2\pi\}.
\] Its normalized \(\mathcal H^1\) measure is \(2\pi\sqrt t\). One can
check this by the injective arc parametrization on \([0,2\pi)\): the
speed is \(\sqrt t\), so its length is
\(\int_0^{2\pi}\sqrt t\,d\theta\). At \(t=0\) the level is the singleton
\(\{(0,0)\}\), whose \(\mathcal H^1\) measure is zero. Outside \([0,1]\)
the level in \(D\) is empty.

The only critical point is \((0,0)\) because \(\nabla g=(2x,2y)\). The
only critical value is therefore \(0\). Here this is a direct
calculation, requiring no appeal to Sard's theorem. The regular levels
are circles; the exceptional level is a point. Treating it separately
avoids a division by zero.

\textbf{2. First apply coarea without dividing.} The weighted coarea
formula for this smooth scalar map is \[
\int_{\mathbb R^2}w(z)|\nabla g(z)|\,d\mathcal L^2(z)
=\int_{\mathbb R}\left(\int_{g^{-1}(\{t\})}w(z)\,d\mathcal H^1(z)\right)dt
\] for nonnegative Borel \(w\), with the usual convention that the
product is zero where the Jacobian is zero. Taking \(w=\mathbf1_D\)
gives \[
\int_D|\nabla g|\,d\mathcal L^2
=\int_0^1 2\pi\sqrt t\,dt
=\frac{4\pi}{3}.
\] The value of the inner integral at \(t=0\) is exactly zero. The value
at the endpoint \(t=1\) does not affect its one-dimensional Lebesgue
integral.

\textbf{3. Divide only on the regular part.} Define a Borel function
\(w:\mathbb R^2\to[0,\infty]\) by \[
w(z)=
\begin{cases}
\displaystyle\frac{\phi(g(z))}{2|z|},&z\in D\setminus\{0\},\\[4pt]
0,&z=0\text{ or }z\notin D.
\end{cases}
\] On \(D\setminus\{0\}\) one has \(wJ_1g=\phi\circ g\). At zero the two
functions may have different values, but this is a \(\mathcal L^2\)-null
set. Consequently their nonnegative Lebesgue integrals agree, even if
\(\phi(0)=+\infty\).

For \(0<t\leq1\), \(w\) is constant on the level circle, with value
\(\phi(t)/(2\sqrt t)\). Hence the inner level integral is \[
\int_{g^{-1}(\{t\})}w\,d\mathcal H^1
=\frac{\phi(t)}{2\sqrt t}\,2\pi\sqrt t
=\pi\phi(t).
\] At \(t=0\), the definition \(w(0)=0\) makes the inner integral zero.
Thus coarea proves \[
\int_D\phi\circ g\,d\mathcal L^2
=\int_{(0,1]}\pi\phi(t)\,dt
=\pi\int_0^1\phi(t)\,dt.
\] Only the integrals ignore the endpoints; no claim that
\(1/|\nabla g|\) is defined at every point has been made.

\textbf{4. Read off useful consequences.} For a real exponent \(a\),
assign any finite value at the origin to the function
\(z\mapsto|z|^{2a}\) if its literal formula is undefined there. On
\((0,1]\) use \(\phi(t)=t^a\), and assign \(\phi(0)=0\). The preceding
identity gives \[
\int_D|z|^{2a}\,d\mathcal L^2(z)
=\begin{cases}
\displaystyle\frac{\pi}{a+1},&a>-1,\\[4pt]
+\infty,&a\leq-1.
\end{cases}
\] In particular \(|z|^{-1}\) is integrable on the disk, with integral
\(2\pi\), whereas \(|z|^{-2}\) is not. The divergence occurs throughout
arbitrarily small neighborhoods of zero; changing the value at zero
cannot repair it.

For \(0\leq a_0<b_0\leq1\), taking \(\phi=\mathbf1_{(a_0,b_0)}\) gives
the annular area \(\pi(b_0-a_0)\) for \(a_0<|z|^2<b_0\). Equivalently,
if a random point is uniform in \(D\) with probability measure
\(\pi^{-1}\mathcal L^2|_D\), its squared radius is uniform on \([0,1]\).
These are both consequences of the same pushforward identity, not
additional assumptions about polar coordinates.

\subsection{Why the theorem applies}\label{why-the-theorem-applies-11}

The locally Lipschitz weighted coarea theorem applies to \(g\) on
\(\mathbb R^2\), with the Borel weights above. Its Jacobian is the norm
of the gradient because the target dimension is one. This polynomial is
smooth and locally Lipschitz; global Lipschitz continuity on all of
\(\mathbb R^2\) is unnecessary. If using a formulation requiring a
globally Lipschitz map, use \(G(z)=\min\{|z|^2,4\}\). The function
\(r\mapsto\min\{r^2,4\}\) is \(4\)-Lipschitz on \([0,\infty)\), and
\(z\mapsto|z|\) is \(1\)-Lipschitz, so \(G\) is globally \(4\)-Lipschitz
and agrees with \(g\) on a neighborhood of \(D\).

For signed Borel \(\phi\) with \(\int_0^1|\phi(t)|\,dt<\infty\), first
apply the nonnegative identity to \(|\phi|\). This proves absolute
integrability on \(D\); applying the formula to \(\phi^+\) and
\(\phi^-\) then gives the signed identity. If both sign parts have
infinite integral, their difference is undefined.

\subsection{Pitfalls}\label{pitfalls-11}

\begin{itemize}
\tightlist
\item
  Coarea gives a factor \(|\nabla g|=2r\), not its reciprocal. A
  reciprocal appears only because we deliberately choose a compensating
  weight on the regular part.
\item
  A critical level need not be a regular curve. Its contribution must be
  checked, rather than parametrized as though its dimension were
  unchanged.
\item
  The critical point being null does not by itself establish
  integrability of a singular weight near it; the power computation
  supplies the needed neighborhood estimate.
\item
  These are scalar integrals over levels. No orientation or vector flux
  occurs.
\end{itemize}

\subsection{Reusable technique}\label{reusable-technique-11}

For an integral depending on a scalar map \(g\), compute the geometry of
\(g^{-1}(\{t\})\) and the coarea Jacobian, isolate the critical set, and
choose a compensating weight only where the Jacobian is nonzero. Tags:
\texttt{coarea}, \texttt{level-\allowbreak{}sets}, \texttt{radial-\allowbreak{}integral},
\texttt{critical-\allowbreak{}value}, \texttt{singular-\allowbreak{}weight},
\texttt{pushforward-\allowbreak{}measure}.

\subsection{Sources}\label{sources-11}

\begin{itemize}
\tightlist
\item
  Leon Simon,
  \href{https://math.stanford.edu/~lms/ntu-gmt-text.pdf}{\emph{Introduction
  to Geometric Measure Theory}}, Chapter 2, §7, Theorem 7.3, Remark
  7.4(1), and Theorem 7.5, printed pp.~76, 80; source ID
  \texttt{simon-\allowbreak{}gmt}. The weighted coarea identity is the theorem used;
  the critical point and all integrals in this example are computed
  directly. Chapter 2, §3, Example 3.5(1) connects normalized
  \(\mathcal H^1\) with arc length.
\end{itemize}

\section{E13 --- A singular map loses all its increase under
differentiation}\label{E13}

\subsection{Problem}\label{problem-12}

Construct a continuous nondecreasing map \(C:[0,1]\to[0,1]\) with
\(C(0)=0\), \(C(1)=1\), and \(C'=0\) almost everywhere. Compare \[
\int_0^1|C'(x)|\,dx
\quad\text{with}\quad
\int_{\mathbb R}\mathcal H^0\bigl(C^{-1}(\{y\})\bigr)\,dy.
\] Explain exactly why the Lipschitz area formula and the fundamental
theorem for absolutely continuous functions cannot be applied to \(C\).
Use the construction to exhibit concentration that prevents exchanging a
derivative limit with an integral, even though the functions themselves
converge uniformly.

\subsection{Data and hypotheses}\label{data-and-hypotheses-12}

The source and target use completed one-dimensional Lebesgue measure
\(\mathcal L^1\). Normalized \(\mathcal H^1\) equals \(\mathcal L^1\) on
a line; the covering normalization in dimension \(k\geq1\) is
\(\alpha_k(\operatorname{diam}U)^k/2^k\). The measure \(\mathcal H^0\)
counts points, so an infinite fibre has value \(+\infty\).

Derivatives will be used as almost-everywhere defined functions. Any
undefined values on an explicitly identified null set may be assigned
zero to obtain a measurable representative. This convention does not
assert differentiability there. All finite-stage approximations below
are piecewise linear and absolutely continuous; their limit will fail
absolute continuity.

\subsection{Solution}\label{solution-12}

\textbf{1. Build the set and the approximations.} Let \(K_0=[0,1]\).
Obtain \(K_{n+1}\) from \(K_n\) by removing the open middle third of
each component interval. The set \(K_n\) has \(2^n\) disjoint closed
component intervals, each of length \(3^{-n}\). Thus \[
\mathcal L^1(K_n)=2^n3^{-n}=\left(\frac23\right)^n.
\] The compact Cantor set \(K=\bigcap_{n=0}^{\infty}K_n\) is null by
continuity of finite measure from above. It therefore has empty
interior, because every nondegenerate interval has positive Lebesgue
measure.

Define \(C_n\) to be the continuous nondecreasing function that is
constant on each gap removed by stage \(n\), rises linearly by
\(2^{-n}\) on each component interval of \(K_n\), and satisfies
\(C_n(0)=0\). These rules uniquely determine it, with \(C_n(1)=1\). For
example \(C_0(x)=x\), while \(C_1\) rises from \(0\) to \(1/2\) on
\([0,1/3]\), is constant \(1/2\) on \([1/3,2/3]\), and rises to \(1\) on
\([2/3,1]\).

All later approximations retain the endpoint values of every stage-\(n\)
component interval. Within such an interval both \(C_m\) and \(C_n\),
for \(m\geq n\), remain between endpoint values separated by \(2^{-n}\);
on previously removed gaps they agree. Therefore \[
\sup_{x\in[0,1]}|C_m(x)-C_n(x)|\leq2^{-n}\qquad(m\geq n).
\] The sequence converges uniformly to a continuous nondecreasing
function \(C\). The endpoint values persist, so \(C(0)=0\), \(C(1)=1\),
and the intermediate value theorem gives \(C([0,1])=[0,1]\). This
construction avoids all ambiguity of ternary expansions at interval
endpoints.

\textbf{2. Compute the almost-everywhere derivative.} Every
\(x\in[0,1]\setminus K\) lies in an open gap removed at a finite stage.
All sufficiently late \(C_n\), and hence \(C\), are constant on that
gap. Thus \(C'(x)=0\) at every such \(x\). Since \(\mathcal L^1(K)=0\),
\(C'=0\) almost everywhere. Assigning zero on \(K\) yields the zero
function as a measurable representative, and therefore \[
\int_0^1|C'|\,d\mathcal L^1=0,
\qquad C(1)-C(0)=1.
\] We do not claim that \(C'\) exists at every point of \(K\).

\textbf{3. Compute the multiplicity on the target.} Let \(P\) be the set
of constant values of \(C\) on its deleted open gaps. There are
countably many gaps, so \(P\) is countable; these plateau values are
dyadic rationals. If two points \(a<b\) satisfy \(C(a)=C(b)=y\),
monotonicity forces \(C\) to be constant on \([a,b]\). Because \(K\) has
empty interior, \((a,b)\) contains part of a deleted gap, so \(y\in P\).
Consequently every \(y\in[0,1]\setminus P\) has exactly one preimage,
using surjectivity for existence. Every \(y\in P\) has infinitely many
preimages, including an entire gap. Thus the counting function \(N\) has
values \[
N(C,[0,1],y)=
\begin{cases}
1,&y\in[0,1]\setminus P,\\
+\infty,&y\in P,\\
0,&y\notin[0,1].
\end{cases}
\] It is measurable, and the infinite values on the countable null set
\(P\) do not change its nonnegative Lebesgue integral. Hence \[
\int_{\mathbb R}N(C,[0,1],y)\,dy=1\ne0=\int_0^1|C'|\,dx.
\] This is a failed extension of the Lipschitz area identity to a
continuous, almost-everywhere differentiable map.

\textbf{4. Identify the exact regularity failure.} On the leftmost
stage-\(n\) interval, \(C(3^{-n})-C(0)=2^{-n}\). The corresponding
difference quotient is \((3/2)^n\), so no finite Lipschitz constant
exists.

Absolute continuity fails directly as well. For the \(2^n\) basic
intervals \([a_{n,j},b_{n,j}]\) at stage \(n\), \[
\sum_{j=1}^{2^n}(b_{n,j}-a_{n,j})=(2/3)^n\longrightarrow0,
\qquad
\sum_{j=1}^{2^n}|C(b_{n,j})-C(a_{n,j})|
=\sum_{j=1}^{2^n}2^{-n}=1.
\] Taking \(\varepsilon=1/2\) contradicts the defining
absolute-continuity condition for every proposed \(\delta>0\).

There is also an image-of-null-set obstruction. Every plateau value has
a preimage at a gap endpoint belonging to \(K\), while any preimage
outside \(K\) lies in a gap. Together with surjectivity this proves
\(C(K)=[0,1]\). Thus this map sends a null set to a set of measure one.
In contrast, a Lipschitz map on a line sends null sets to null sets:
cover the source null set by intervals of arbitrarily small total
length, then cover their images by intervals whose lengths are at most
the Lipschitz constant times those lengths. Discarding a source null set
is therefore unsafe for this singular map.

\textbf{5. See concentration in the derivatives.} At all but finitely
many breakpoints, \[
C_n'(x)=\left(\frac32\right)^n\mathbf1_{K_n}(x).
\] Choose the displayed right side as an almost-everywhere
representative, denoted \(q_n\). Then \[
\int_0^1q_n(x)\,dx
=\left(\frac32\right)^n\left(\frac23\right)^n=1,
\qquad q_n\longrightarrow0\quad\mathcal L^1\text{-almost everywhere}.
\] Thus \(\int\lim q_n=0\) while \(\lim\int q_n=1\). The mass stays
inside a bounded interval and concentrates on the shrinking sets \(K_n\)
containing the null set \(K\); it does not escape to infinity.

There is no common integrable dominator. If \(q_n\leq G\) almost
everywhere for \(G\in L^1\), then \(\int_{K_n}|G|\geq1\) for every
\(n\). Dominated convergence applied to \(|G|\mathbf1_{K_n}\to0\) almost
everywhere would instead give \(\int_{K_n}|G|\to0\), a contradiction.
The sequence also fails uniform integrability, since
\(\mathcal L^1(K_n)\to0\) but \(\int_{K_n}q_n=1\).

\subsection{Why the theorem applies}\label{why-the-theorem-applies-12}

The continuity-of-measure, uniform-limit, and nonnegative integration
arguments used in the construction have their stated hypotheses. The
area formula is used only as a comparison theorem: its Lipschitz
hypothesis fails for \(C\), so its equality is not asserted. The
absolute-continuity fundamental theorem similarly does not apply to
\(C\).

A useful positive replacement is the following. If \(F_n\) are
absolutely continuous on \([0,1]\), \(F_n(0)\to a\), \(F_n'\to q\)
almost everywhere, and \(|F_n'|\leq G\in L^1\) almost everywhere, then
dominated convergence gives \[
\sup_{x\in[0,1]}\left|F_n(x)-a-\int_0^xq(t)\,dt\right|
\leq |F_n(0)-a|+\int_0^1|F_n'-q|\,dt\longrightarrow0.
\] Here \(q\in L^1\) because \(|q|\leq G\) almost everywhere, and the
limit is absolutely continuous with derivative \(q\) almost everywhere.
This is a precise sufficient condition for retaining derivatives under
the limit; the Cantor approximations violate the domination condition.

\subsection{Pitfalls}\label{pitfalls-12}

\begin{itemize}
\tightlist
\item
  Uniform convergence of functions does not supply domination of their
  derivatives.
\item
  An almost-everywhere zero derivative alone does not make a continuous
  function constant. Absolute continuity restores that conclusion.
\item
  An infinite multiplicity at countably many target values is harmless
  for this target integral; a null source set with a full-measure image
  is a different issue.
\item
  This does not contradict a smooth Sard theorem: \(C\) is not a smooth
  map, and the exceptional domain set cannot be removed while preserving
  its image measure.
\end{itemize}

\subsection{Reusable technique}\label{reusable-technique-12}

Track where derivative mass lives, not just the pointwise limit. Check
absolute continuity or the actual regularity required by
area/change-of-variables formulas; test whether null source sets can
have large images. Tags: \texttt{singular-\allowbreak{}map},
\texttt{cantor-\allowbreak{}function}, \texttt{absolute-\allowbreak{}continuity},
\texttt{area-\allowbreak{}formula-\allowbreak{}failure}, \texttt{concentration},
\texttt{dominated-\allowbreak{}convergence-\allowbreak{}failure}, \texttt{uniform-\allowbreak{}integrability},
\texttt{derivative-\allowbreak{}limit}.

\subsection{Sources}\label{sources-12}

\begin{itemize}
\tightlist
\item
  John K. Hunter,
  \href{https://www.math.ucdavis.edu/~hunter/measure_theory/measure_notes.pdf}{\emph{Measure
  Theory}}, Example 2.14 (Cantor-set nullity), Example 2.37 (singular
  Cantor measure), and the opening of Chapter 6, printed pp.~19, 31, 63;
  source ID \texttt{hunter-\allowbreak{}measure}. These give background on the
  standard set and the absolute-continuity obstruction; the construction
  and calculations above are supplied explicitly.
\item
  Leon Simon,
  \href{https://math.stanford.edu/~lms/ntu-gmt-text.pdf}{\emph{Introduction
  to Geometric Measure Theory}}, Chapter 2, §3, Theorem 3.3 and Remark
  3.4(2), printed pp.~57--58; source ID \texttt{simon-\allowbreak{}gmt}. The
  Lipschitz area theorem is the precise formula whose unjustified
  extension is disproved here.
\end{itemize}

\section{E14 --- Lost integral mass: concentration, escape, and
dilution}\label{E14}

\subsection{Problem}\label{problem-13}

Compare the limits and integrals of these measurable functions: \[
f_n(x)=n\mathbf1_{(0,1/n)}(x)\quad(x\in(0,1)),\qquad
h_n(x)=\mathbf1_{(n,n+1)}(x)\quad(x\in\mathbb R).
\] Explain why dominated convergence fails. Identify sufficient
replacement hypotheses, keeping finite and infinite measure spaces
distinct.

\subsection{Data and hypotheses}\label{data-and-hypotheses-13}

All measures are one-dimensional Lebesgue measure. The symbol \(f_n\)
denotes a function; \(f_n(x)\) denotes its value. All three sequences
used below consist of nonnegative Borel functions. Convergence in
measure on a measure space means \(\mu(\{|u_n-u|>\varepsilon\})\to0\)
for each \(\varepsilon>0\). ``Locally in measure'' on \(\mathbb R\)
means this on every bounded interval.

We use \textbf{value-tail uniform integrability} to mean \[
\lim_{K\to\infty}\sup_n\int_{\{|u_n|>K\}}|u_n|\,d\mu=0.
\] On finite measure spaces this is the usual uniform-integrability
criterion. On infinite measure spaces some definitions additionally
require spatial tightness; here that condition will always be stated
separately.

\subsection{Solution}\label{solution-13}

\textbf{1. Concentration on a shrinking set.} For each fixed
\(x\in(0,1)\), eventually \(n>1/x\), so \(f_n(x)=0\). Nevertheless \[
\int_0^1 f_n(x)\,dx=n\lambda^1((0,1/n))=1.
\] Thus \(f_n\to0\) pointwise everywhere on its domain, while
\(\int f_n\not\to\int0\). It also converges in measure since, for
\(n>\varepsilon\), \[
\lambda^1(\{f_n>\varepsilon\})=1/n\longrightarrow0.
\] There is no common integrable majorant \(G\): if \(G\ge f_n\) almost
everywhere for all \(n\), take a common full-measure set by removing the
countable union of exceptional sets. Then \(\int_0^{1/n}G\ge1\) for
every \(n\), contradicting absolute continuity of the integral of an
\(L^1\) function as \(1/n\to0\). Its value-tail integrals also show
precisely what fails: for any \(K\), choose an integer \(n>K\) to obtain
\[
\int_{\{f_n>K\}}f_n\,dx=1.
\]

\textbf{2. Escape to infinity without large values.} For each
\(x\in\mathbb R\), \(h_n(x)=0\) eventually, but
\(\int_{\mathbb R}h_n=1\). The functions are bounded by \(1\), which is
not integrable on \(\mathbb R\). A common \(L^1\) majorant is impossible
because the intervals \((n,n+1)\) are disjoint and \[
\int_{\mathbb R}G\,dx\ge\sum_{n=1}^{\infty}\int_n^{n+1}G\,dx\ge\sum_{n=1}^{\infty}1=\infty.
\] Here value-tail uniform integrability holds: for \(K\ge1\) every tail
integral is zero. What fails is \textbf{spatial tightness}: \[
\lim_{R\to\infty}\sup_n\int_{\{|x|>R\}}|h_n(x)|\,dx=1,
\] not zero. Also \(h_n\to0\) locally in measure but not in measure on
all of \(\mathbb R\), since \(\lambda^1(\{h_n>1/2\})=1\).

\textbf{3. Even global convergence in measure is insufficient on an
infinite space.} Set \(r_n(x)=n^{-1}\mathbf1_{(0,n)}(x)\) on
\(\mathbb R\). Then \(r_n\to0\) pointwise and uniformly, hence in
measure, and \(\int r_n=1\). Its values are at most \(1\), so it is
value-tail uniformly integrable. But for fixed \(R>0\), \[
\int_{\{|x|>R\}}r_n\,dx=\max(1-R/n,0),
\] whose supremum over \(n\) is \(1\). Mass spreads out at smaller
heights. This is why a finite-measure version of a theorem cannot
silently be applied to \(\mathbb R\).

\subsection{Why the theorem applies}\label{why-the-theorem-applies-13}

Useful sufficient conditions are:

\begin{itemize}
\tightlist
\item
  On any measure space, if \(u_n\to u\) almost everywhere and
  \(|u_n|\le G\) almost everywhere for one \(G\in L^1\), dominated
  convergence gives \(\|u_n-u\|_1\to0\) and therefore convergence of
  integrals.
\item
  On a \textbf{finite} measure space, if \(u_n\) and \(u\) are real
  measurable functions, \(u_n\in L^1\), \(u_n\to u\) in measure, and the
  sequence is value-tail uniformly integrable, then \(u\in L^1\) and
  \(\|u_n-u\|_1\to0\).
\item
  On \(\mathbb R\), for real measurable \(u_n,u\) with \(u_n\in L^1\),
  the same conclusion follows from local convergence in measure,
  value-tail uniform integrability, and spatial tightness
  \(\lim_{R\to\infty}\sup_n\int_{|x|>R}|u_n|=0\).
\end{itemize}

For completeness, the finite-measure assertion follows by truncation:
let \(T_K(s)=\max(-K,\min(s,K))\) for real functions. On a finite space,
convergence in measure gives \[
\int|T_K(u_n)-T_K(u)|\le\varepsilon\mu(X)+2K\mu(\{|u_n-u|>\varepsilon\})\longrightarrow\text{at most }\varepsilon\mu(X).
\] Uniform integrability controls \(\|u_n-T_K(u_n)\|_1\) uniformly. A
subsequence converges almost everywhere (choose error probabilities
summable); Fatou bounds the corresponding tail for \(u\) by the same
uniform tail bound. Let \(K\to\infty\) and then \(\varepsilon\to0\). On
\(\mathbb R\), apply this finite-space argument on \([-R,R]\) and
control its complement by spatial tightness. To bound the tail of \(u\),
choose a diagonal subsequence converging almost everywhere on each
\([-m,m]\): at its \(j\)th step make the measure of the error set
\(\{|u_{n_j}-u|>1/j\}\cap[-j,j]\) less than \(2^{-j}\). The
summable-error argument gives almost-everywhere convergence on all of
\(\mathbb R\), and Fatou then bounds \(\int_{|x|>R}|u|\) by the uniform
tail bound. These are analytic hypotheses, not conclusions inferred from
numerical samples.

\subsection{Pitfalls}\label{pitfalls-13}

Each \(\int u_n\) being finite, or their values being uniformly bounded,
is insufficient. Uniform convergence alone does not control integrals
over an infinite-measure domain. No example contradicts Fatou: it only
asserts \(0=\int\liminf u_n\le\liminf\int u_n=1\) here.

\subsection{Reusable technique}\label{reusable-technique-13}

Before moving a limit inside an integral, inspect large values on small
sets and mass outside large bounded sets separately. Tags:
\texttt{dominated-\allowbreak{}convergence}, \texttt{uniform-\allowbreak{}integrability},
\texttt{concentration}, \texttt{escape-\allowbreak{}to-\allowbreak{}infinity}, \texttt{tightness},
\texttt{counterexample}.

\subsection{Sources}\label{sources-13}

\begin{itemize}
\tightlist
\item
  John K. Hunter,
  \href{https://www.math.ucdavis.edu/~hunter/measure_theory/measure_notes.pdf}{\emph{Measure
  Theory}}, Examples 4.18--4.20 (the standard concentration, dilution
  and escape families), Theorem 4.24 (dominated convergence), the
  following discussion of absolute continuity and uniform integrability,
  and Chapter 4's Fatou lemma. The computations and additional
  uniform-integrability/tightness analysis here are worked through
  explicitly. Source ID: \texttt{hunter-\allowbreak{}measure}.
\item
  The uniform-integrability implications used here are proved in the
  paragraph above; the source is not being cited as containing this
  exact formulation or atlas.
\end{itemize}

\section{E15 --- A parameter derivative misses a concentrating boundary
layer}\label{E15}

\subsection{Problem}\label{problem-14}

On \(x\in(0,1)\), for \(t\ge0\), define \[
f(t,x)=\begin{cases}\dfrac{t^2}{x^2+t^2},&t>0,\\0,&t=0,\end{cases}
\qquad I(t)=\int_0^1f(t,x)\,dx.
\] Compare \(I'_+(0)\) with \(\int_0^1\partial_t^+f(0,x)\,dx\). Locate
the precise missing hypothesis for differentiation under the integral
sign.

\subsection{Data and hypotheses}\label{data-and-hypotheses-14}

The measure is Lebesgue measure \(\lambda^1\) on \((0,1)\). For every
\(t\), \(x\mapsto f(t,x)\) is Borel and satisfies \(0\le f(t,x)\le1\).
For each fixed \(x>0\), \(t\mapsto f(t,x)\) has right derivative at
zero. We ask about a derivative of the \textbf{integral-valued function}
\(I\), not about a derivative of one fixed value \(f(t,x)\).

\subsection{Solution}\label{solution-14}

\textbf{1. Evaluate the integral before differentiating.} For \(t>0\),
substitute \(x=tu\): \[
I(t)=t\int_0^{1/t}\frac{1}{1+u^2}\,du=t\arctan(1/t),\qquad I(0)=0.
\] It follows that \[
I'_+(0)=\lim_{t\downarrow0}\frac{I(t)-I(0)}{t}
=\lim_{t\downarrow0}\arctan(1/t)=\frac\pi2.
\]

\textbf{2. Differentiate each fixed section.} For every fixed \(x>0\),
\[
\partial_t^+f(0,x)=\lim_{t\downarrow0}\frac{f(t,x)-f(0,x)}{t}
=\lim_{t\downarrow0}\frac{t}{x^2+t^2}=0.
\] Thus \(\int_0^1\partial_t^+f(0,x)\,dx=0\), which differs from
\(I'_+(0)\).

\textbf{3. Find the boundary layer.} The difference-quotient functions
are \[
q_t(x)=\frac{t}{x^2+t^2}\ge0.
\] They tend to zero at each fixed \(x>0\), but \[
\int_0^1q_t(x)\,dx=\arctan(1/t)\longrightarrow\pi/2.
\] For \(0<t\le1\) the interval \((0,t)\) already carries a fixed amount
of mass: \[
\int_0^t q_t(x)\,dx=\int_0^1\frac{du}{1+u^2}=\pi/4.
\] This proves that the family of quotients has no common \(L^1(0,1)\)
majorant on any neighborhood \(0<t<\delta\): otherwise absolute
continuity of its integral would force the displayed integrals to tend
to zero.

The derivative family also has no local common integrable majorant. For
\(t>0\), \[
\partial_t f(t,x)=\frac{2tx^2}{(x^2+t^2)^2}.
\] To avoid any uncountable exceptional-set issue, suppose \(G\in L^1\)
bounds this derivative for almost every \(x\) for \textbf{each}
\(t\in(0,\delta)\). Take \(t_j\downarrow0\) with \(2t_j<1\). On
\((t_j,2t_j)\), put \(x=t_ju\) to obtain \[
\int_{t_j}^{2t_j}\partial_t f(t_j,x)\,dx
=\int_1^2\frac{2u^2}{(1+u^2)^2}\,du>0,
\] independent of \(j\). Since \(G\) dominates the integrand at these
countably many parameters, absolute continuity again gives a
contradiction. This even rules out bounds whose null exceptional set was
initially allowed to depend on the parameter.

\subsection{Why the theorem applies}\label{why-the-theorem-applies-14}

Ordinary dominated convergence \textbf{does} give \(I(t)\to0\):
\(f(t,x)\to0\) and \(|f(t,x)|\le1\in L^1(0,1)\). It says nothing about
the derivative unless the quotients are controlled.

A precise sufficient differentiation rule is as follows. Let
\((X,\mathcal A,\mu)\) be a measure space and \(J\) an open real
interval containing \(t_0\). Assume \(F(t,\cdot)\) is measurable for
every \(t\in J\), \(F(t_0,\cdot)\in L^1(\mu)\), and there is a
\textbf{common} full-measure set on which \(t\mapsto F(t,x)\) is
differentiable throughout \(J\) and \[
|\partial_tF(t,x)|\le G(x)\quad(t\in J),\qquad G\in L^1(\mu).
\] Then \(F(t,\cdot)\in L^1\) near \(t_0\) and \[
\left.\frac{d}{dt}\right|_{t=t_0}\int_XF(t,x)\,d\mu(x)
=\int_X\partial_tF(t_0,x)\,d\mu(x).
\] Indeed the mean value theorem bounds the difference quotients by
\(G\), and dominated convergence proves the identity along every
sequence of nonzero increments tending to zero. The same proof works for
a right derivative at a left endpoint, using a half-interval and right
differentiability at the endpoint. In this example the required bound
fails exactly where \(x\) and \(t\) approach zero together.

\subsection{Pitfalls}\label{pitfalls-14}

\begin{itemize}
\tightlist
\item
  An integrable bound on \(f\) is not an integrable bound on its
  parameter derivative or difference quotients.
\item
  Differentiability at every fixed \(x\) does not imply differentiation
  can pass through integration.
\item
  A related warning for a different exchange is \(v_n(x)=\sin(nx)/n\) on
  \([-1,1]\): \(v_n\to0\) uniformly, but \(v_n'(0)=1\) for all \(n\),
  whereas the limit function has derivative \(0\). Uniform convergence
  of functions alone does not justify exchanging a limit and a
  derivative.
\end{itemize}

\subsection{Reusable technique}\label{reusable-technique-14}

Write the actual difference quotients and search for a common integrable
bound. When the quotient concentrates, integrate it over a
parameter-dependent shrinking interval to prove that no such bound
exists. Tags: \texttt{differentiate-\allowbreak{}under-\allowbreak{}integral},
\texttt{boundary-\allowbreak{}layer}, \texttt{concentration},
\texttt{dominated-\allowbreak{}convergence}, \texttt{counterexample}.

\subsection{Sources}\label{sources-14}

\begin{itemize}
\tightlist
\item
  John K. Hunter,
  \href{https://www.math.ucdavis.edu/~hunter/measure_theory/measure_notes.pdf}{\emph{Measure
  Theory}}, Theorem 4.24 (the DCT used in the proof). Source ID:
  \texttt{hunter-\allowbreak{}measure}.
\item
  Keith Conrad,
  \href{https://kconrad.math.uconn.edu/blurbs/analysis/diffunderint.pdf}{\emph{Differentiating
  under the integral sign}}, Section 12 and Theorem 12.3, for classical
  hypotheses and a warning example. The measurable sufficient rule here
  is explicitly proved above; the present boundary-layer example is
  independently worked. Source ID: \texttt{conrad-\allowbreak{}diff}.
\end{itemize}

\section{E16 --- An exponential kernel retrieves a boundary
value}\label{E16}

\subsection{Problem}\label{problem-15}

Let \(g:[0,\infty)\to\mathbb R\) be bounded and Borel measurable, and
suppose it is right-continuous at \(0\). Prove \[
\lim_{n\to\infty}\int_0^{\infty}ne^{-nx}g(x)\,dx=g(0).
\] Give a direct error estimate, and identify why applying dominated
convergence in the original variable would be unsafe.

\subsection{Data and hypotheses}\label{data-and-hypotheses-15}

The measure is one-dimensional Lebesgue measure on \([0,\infty)\).
Choose \(M\ge0\) with \(|g(x)|\le M\) for all \(x\ge0\). The kernel
function is \(k_n:x\mapsto ne^{-nx}\); its value at \(x\) is \(k_n(x)\).
It satisfies \[
k_n\ge0,\qquad\int_0^{\infty}k_n(x)\,dx=1,\qquad
\int_\delta^{\infty}k_n(x)\,dx=e^{-n\delta}\quad(\delta>0).
\] Every displayed integral of \(k_ng\) is absolutely convergent because
\(\int|k_ng|\le M\). The value \(g(0)\) is recoverable under the stated
continuity assumption; Lebesgue integration by itself ignores a single
point.

\subsection{Solution}\label{solution-15}

\textbf{1. Rescale the moving kernel.} The substitution \(u=nx\), with
\(dx=du/n\), gives \[
\int_0^{\infty}ne^{-nx}g(x)\,dx
=\int_0^{\infty}e^{-u}g(u/n)\,du.
\] For each fixed \(u\ge0\), right-continuity gives \(g(u/n)\to g(0)\).
Moreover \[
|e^{-u}g(u/n)|\le Me^{-u},\qquad
\int_0^{\infty}Me^{-u}\,du=M<\infty.
\] Dominated convergence now applies on the fixed \(u\)-domain and gives
\[
\lim_n\int_0^{\infty}e^{-u}g(u/n)\,du
=g(0)\int_0^{\infty}e^{-u}\,du=g(0).
\]

\textbf{2. Obtain a direct epsilon and tail proof.} For any
\(\varepsilon>0\), choose \(\delta>0\) such that
\(|g(x)-g(0)|<\varepsilon\) when \(0\le x<\delta\). Normalization of
\(k_n\) yields \[
\begin{aligned}
\left|\int_0^{\infty}k_n(x)g(x)\,dx-g(0)\right|
&\le\int_0^\delta k_n(x)|g(x)-g(0)|\,dx
 +\int_\delta^{\infty}k_n(x)|g(x)-g(0)|\,dx\\
&\le\varepsilon+2M e^{-n\delta}.
\end{aligned}
\] For this fixed \(\delta\), the tail tends to zero. Then let
\(\varepsilon\downarrow0\).

\textbf{3. A useful convergence rate.} If, additionally,
\(|g(x)-g(0)|\le Lx\) for every \(x\ge0\) (in particular if \(g\) is
globally \(L\)-Lipschitz), then \[
\left|\int_0^{\infty}k_n(x)g(x)\,dx-g(0)\right|
\le L\int_0^{\infty}ne^{-nx}x\,dx
=\frac Ln\int_0^{\infty}ue^{-u}\,du=\frac Ln.
\] The last integral is \(1\) by integration by parts. If only a local
Lipschitz bound holds for \(0\le x<\delta\), the same split gives the
valid bound \(L/n+2Me^{-n\delta}\).

\subsection{Why the theorem applies}\label{why-the-theorem-applies-15}

The rescaled integrands \(u\mapsto e^{-u}g(u/n)\) are Borel, converge
pointwise, and have the explicit common \(L^1\) majorant \(Me^{-u}\).
These are precisely the dominated-convergence hypotheses. The
substitution is a positive linear change of variables, so its Jacobian
is \(1/n\); it is valid for these absolutely integrable measurable
functions.

In the original variable \(k_n(x)\to0\) for every \(x>0\), while
\(\int k_n=1\). Hence no common \(L^1\) function dominates all \(k_n\).
More directly, if \(G\ge k_n\) almost everywhere for all \(n\), then \[
\int_0^{1/n}G\,dx\ge\int_0^{1/n}ne^{-nx}\,dx=1-e^{-1},
\] contradicting absolute continuity of an \(L^1\) integral. Rescaling
resolves the concentration instead of incorrectly ignoring it. In
measure language the probability measures \(k_n(x)\,dx\) converge weakly
against bounded continuous test functions on \([0,\infty)\) to the unit
point mass at \(0\); they do not converge to the zero measure merely
because their densities tend to zero almost everywhere.

\subsection{Pitfalls}\label{pitfalls-15}

\begin{itemize}
\tightlist
\item
  Do not infer the limit integral is zero from \(ne^{-nx}g(x)\to0\) for
  \(x>0\): a common integrable bound has not been supplied.
\item
  If right-continuity is removed, the answer need not equal the chosen
  value at zero. For \(g(0)=1\) and \(g(x)=0\) for \(x>0\), every
  integral is zero.
\item
  Boundedness and continuity at zero supply different parts of the
  proof: boundedness controls the far tail, while continuity identifies
  the limit near the boundary.
\end{itemize}

\subsection{Reusable technique}\label{reusable-technique-15}

For a normalized concentrating kernel, either rescale to a fixed kernel
and use DCT, or split into a small neighborhood and an exponentially
small tail. Tags: \texttt{approximate-\allowbreak{}identity},
\texttt{dominated-\allowbreak{}convergence}, \texttt{rescaling},
\texttt{boundary-\allowbreak{}value}, \texttt{error-\allowbreak{}bound}, \texttt{laplace-\allowbreak{}kernel}.

\subsection{Sources}\label{sources-15}

\begin{itemize}
\tightlist
\item
  John K. Hunter,
  \href{https://www.math.ucdavis.edu/~hunter/measure_theory/measure_notes.pdf}{\emph{Measure
  Theory}}, Theorem 4.24 (dominated convergence) and Chapter 4's
  absolute continuity of the integral. Source ID:
  \texttt{hunter-\allowbreak{}measure}.
\item
  The exponential-kernel calculation, estimates, and rate are proved
  directly above; no distributional formalism is required.
\end{itemize}

\section{A small atlas of failed exchanges}\label{atlas}

This is a retrieval map to complete proofs, not a replacement for them.
It gives outcomes; use \hyperref[problems]{PROBLEMS.md} for the
problem-first route.

\begin{longtable}[]{@{}
  >{\raggedright\arraybackslash}p{(\columnwidth - 4\tabcolsep) * \real{0.2800}}
  >{\raggedright\arraybackslash}p{(\columnwidth - 4\tabcolsep) * \real{0.3600}}
  >{\raggedright\arraybackslash}p{(\columnwidth - 4\tabcolsep) * \real{0.3600}}@{}}
\toprule\noalign{}
\begin{minipage}[b]{\linewidth}\raggedright
Exchange or tempting inference
\end{minipage} & \begin{minipage}[b]{\linewidth}\raggedright
Example and failure mechanism
\end{minipage} & \begin{minipage}[b]{\linewidth}\raggedright
Positive sufficient condition
\end{minipage} \\
\midrule\noalign{}
\endhead
\bottomrule\noalign{}
\endlastfoot
Integrate every section of a Lebesgue-measurable set & E03: a
nonmeasurable set can sit in one exceptional section of a planar null
set. & Borel measurability gives Borel sections everywhere; completed
Lebesgue measurability gives measurable sections almost everywhere. \\
Swap signed iterated integrals & E04: both signed parts have infinite
integral near the origin. E05: the signed matrix has infinite absolute
sum. & Joint measurability and finite integral of the absolute value
give Fubini; nonnegativity gives Tonelli with extended values. \\
Pass an almost-everywhere limit through an integral & E14: mass
concentrates at a shrinking interval; its height grows. & DCT with one
integrable majorant; on a finite measure space, convergence in measure
plus uniform integrability gives convergence in \(L^1\). \\
Treat bounded integral mass as sufficient on an infinite space & E14: a
unit bump moves to infinity. Value-tail uniform integrability alone
permits this. & Add spatial tightness to value-tail uniform
integrability, together with convergence in measure on bounded sets (and
an integrable limit). Precise versions are stated in E14. \\
Differentiate a parameter integral at an endpoint & E15:
difference-quotient mass concentrates near zero; pointwise derivatives
miss that mass. & Measurability, pointwise differentiability almost
everywhere, one integrable local bound on the parameter derivative, and
one integrable parameter slice. See the exact interval and endpoint
conditions in E15. \\
Differentiate a uniform limit of differentiable functions & E15 related
example: \(\sin(nx)/n\) tends uniformly to zero while its derivative is
\(1\) at zero for every \(n\). & On a compact interval, uniform
convergence of the derivatives and convergence of the function values at
one point imply a differentiable limit with the expected derivative. The
mean value theorem proves this for differentiable functions; for \(C^1\)
functions one can also use the fundamental theorem of calculus. \\
Replace geometric area by parameter area & E08/E11: repeated or folded
parametrizations count the image more than once. & Injectivity outside
sets whose contributions are zero, or explicit multiplicity in the area
formula. \\
Divide by a coarea Jacobian everywhere & E12: the gradient vanishes at
the origin. & Restrict division to the positive-Jacobian set, then prove
the zero set contributes zero or account for its contribution
separately. \\
Use the area formula for a merely continuous monotone map & E13: the
Cantor staircase has derivative zero almost everywhere but maps onto an
interval. & Lipschitz maps satisfy the area formula; in one dimension
absolute continuity gives the fundamental theorem and the relevant
change-of-variable identities. \\
\end{longtable}

Concentration is also a tool: E16 uses an approximate identity. There
the integration variable is rescaled and a fixed integrable majorant
becomes available. The conclusion is a boundary evaluation, with all its
hypotheses stated, rather than an exchange based solely on pointwise
convergence in the original variable.

\section{Sources and conventions audit}\label{bibliography}

All links below were opened and checked on \textbf{2026-10-15}. These
are author-hosted teaching texts, used for theorem statements and
conventions. The worked examples in this directory are original
expositions; the linked texts' proofs and problem collections are not
reproduced. PDF page numbers below count the cover as page 1. A PDF
viewer may display a different logical page label. Online files can
change.

\subsection{\texorpdfstring{\texttt{hunter-\allowbreak{}measure}}{hunter-measure}}\label{hunter-measure}

John K. Hunter, \emph{Measure Theory}, University of California, Davis,
2011.
\href{https://www.math.ucdavis.edu/~hunter/measure_theory/measure_notes.pdf}{Author's
PDF}.

\begin{itemize}
\tightlist
\item
  Monotone convergence: Theorem 4.6, printed p.41, PDF p.45.
\item
  Dominated convergence: Theorem 4.24, printed p.49,
  \href{https://www.math.ucdavis.edu/~hunter/measure_theory/measure_notes.pdf\#page=53}{PDF
  p.53}; a.e. versions are discussed in §4.6, printed p.50.
\item
  Product-measurable sections: §5.1, Proposition 5.2, printed p.55, PDF
  p.59; function sections: Theorem 5.15, printed p.60.
\item
  Fubini with absolute integrability: Theorem 5.18, printed p.61,
  \href{https://www.math.ucdavis.edu/~hunter/measure_theory/measure_notes.pdf\#page=65}{PDF
  p.65}.
\item
  Completion and exceptional sections: Example 5.20 and Theorem 5.21,
  printed p.62,
  \href{https://www.math.ucdavis.edu/~hunter/measure_theory/measure_notes.pdf\#page=66}{PDF
  p.66}.
\end{itemize}

The completed Lebesgue product statement guarantees measurable and
integrable sections almost everywhere. It does not guarantee
measurability of every section. Theorem 5.15's sentence defining complex
section integrals suppresses existence hypotheses; use the explicit
integrability hypotheses of 5.18/5.21.

\subsection{\texorpdfstring{\texttt{axler-\allowbreak{}mira}}{axler-mira}}\label{axler-mira}

Sheldon Axler, \emph{Measure, Integration \& Real Analysis}, Graduate
Texts in Mathematics 282, Springer, 2020; author's online revision dated
12 June 2026. \href{https://measure.axler.net/MIRA.pdf}{Author's PDF}.

\begin{itemize}
\tightlist
\item
  Tonelli: §5B, Theorem 5.28, printed p.129,
  \href{https://measure.axler.net/MIRA.pdf\#page=144}{PDF p.144}.
\item
  Nonnegative double sums: Corollary 5.31, printed p.131, PDF p.146.
\item
  Fubini: Theorem 5.32, printed p.132,
  \href{https://measure.axler.net/MIRA.pdf\#page=147}{PDF p.147}, with
  exceptional-parameter conventions immediately above and in the proof.
\end{itemize}

These product theorems assume sigma-finite factors. Tonelli accepts
nonnegative product-measurable functions and permits the value infinity;
Fubini requires a finite integral of the absolute value. This is an
independent reference for the distinction used in the examples. For
completed Lebesgue measurability rather than product measurability, also
use \texttt{hunter-\allowbreak{}measure} §5.7.

\subsection{\texorpdfstring{\texttt{simon-\allowbreak{}gmt}}{simon-gmt}}\label{simon-gmt}

Leon Simon, \emph{Introduction to Geometric Measure Theory}, Stanford
University, 2017/2018 NTU lecture draft; last-update line: 7 March 2018.
\href{https://math.stanford.edu/~lms/ntu-gmt-text.pdf}{Author's PDF}.

\begin{itemize}
\tightlist
\item
  Normalization: Chapter 1, definitions 2.1--2.2, printed p.11,
  \href{https://math.stanford.edu/~lms/ntu-gmt-text.pdf\#page=9}{PDF
  p.9}; equality with Lebesgue outer measure: Theorem 2.6, printed p.12,
  PDF p.10.
\item
  Area and multiplicity: Chapter 2, Theorem 3.3, printed p.57,
  \href{https://math.stanford.edu/~lms/ntu-gmt-text.pdf\#page=32}{PDF
  p.32}; weighted fiber sum: Remark 3.4(2), printed p.58, PDF p.33;
  curves and submanifolds: Examples 3.5, printed pp.60--61, PDF
  pp.34--35.
\item
  Coarea: Chapter 2, Theorem 7.3 and weighted Remark 7.4(1), printed
  p.76,
  \href{https://math.stanford.edu/~lms/ntu-gmt-text.pdf\#page=42}{PDF
  p.42}.
\item
  Critical parts of fibers: Theorem 7.5, printed p.80, PDF p.44; the
  stronger smooth Sard conclusion is distinguished in Remark 7.6.
\end{itemize}

The PDF has two printed pages per sheet. This version uses the factor
\texttt{omega\_k\ /\ 2\^{}k} multiplying sums of diameter powers, where
\texttt{omega\_k} is unit-ball volume. Consequently normalized Hausdorff
measure agrees with Lebesgue measure on k-planes. Counting measure is
the k=0 convention.

\subsection{\texorpdfstring{\texttt{hajlasz-\allowbreak{}geometry}}{hajlasz-geometry}}\label{hajlasz-geometry}

Piotr Hajłasz, \emph{Geometric Analysis}, University of Pittsburgh,
undated author-hosted lecture notes.
\href{https://www.pitt.edu/~hajlasz/Notatki/Analysis_4.pdf}{Author's
PDF}.

\begin{itemize}
\tightlist
\item
  Normalization: §3.1, printed/PDF p.19; equality with Lebesgue measure:
  Theorem 3.2, p.20.
\item
  Surface integration: Theorem 3.17 and following discussion, p.25.
\item
  Area and multiplicity with the Lusin N condition: Theorem 3.18 and
  Remark 3.19, pp.26--27.
\item
  Coarea: §3.4, Theorem 3.23, pp.28--29; radial slicing and
  projection/Fubini applications, p.29. Almost-everywhere
  rectifiability: Theorem 3.24, p.29.
\end{itemize}

Normalization agrees with \texttt{simon-\allowbreak{}gmt}. Two evident dimensional
misprints occur in §3.4: the coarea Jacobian should vanish when the
derivative has rank less than \textbf{m}, not n; the area half of
Theorem 3.23 needs outer measure \textbf{H\^{}n}, not H\^{}m. The coarea
half uses the correct fiber and target dimensions. Our area statements
use Simon's 2018 Theorem 3.3 and Remark 3.4(2), whose dimensions are
correct.

\subsection{\texorpdfstring{\texttt{hajlasz-\allowbreak{}analysis}}{hajlasz-analysis}}\label{hajlasz-analysis}

Piotr Hajłasz, \emph{Analysis I}, University of Pittsburgh, undated
author-hosted lecture notes.
\href{https://www.pitt.edu/~hajlasz/Notatki/Analysis\%20I.pdf}{Author's
PDF}.

\begin{itemize}
\tightlist
\item
  Theorem 103 and the Cantor staircase discussion, printed/PDF p.76: a
  continuous increasing function can have derivative zero almost
  everywhere while its endpoint increment is positive.
\item
  Absolute continuity and the fundamental theorem: §8.1, Theorem 106,
  printed/PDF pp.77--78.
\item
  Convolution in L1: Theorem 87, pp.66--67; smooth approximate
  identities: Theorem 92, pp.70--71.
\end{itemize}

The Cantor paragraph prints \texttt{f\textquotesingle{}(0)\ =\ 0\ a.e.};
the variable should be x. The example in this library proves its needed
staircase properties directly. Derivative values on its exceptional null
set may be assigned arbitrarily for Lebesgue integration, but those
assignments do not make the staircase absolutely continuous or
Lipschitz.

\subsection{\texorpdfstring{\texttt{conrad-\allowbreak{}diff}}{conrad-diff}}\label{conrad-diff}

Keith Conrad, \emph{Differentiating Under the Integral Sign}, University
of Connecticut, undated author-hosted expository notes.
\href{https://kconrad.math.uconn.edu/blurbs/analysis/diffunderint.pdf}{Author's
PDF}.

\begin{itemize}
\tightlist
\item
  A sufficient differentiation theorem: §12, Theorem 12.3, printed/PDF
  p.16.
\item
  Local parameter bounds illustrating the theorem: Table 2, p.16, and
  discussion on p.17.
\end{itemize}

Theorem 12.3 treats continuous integrands and parameter derivatives,
dominated by parameter-independent integrable majorants near the
parameter in question. Its infinite-interval integrals are stated as
improper integrals. The library supplies Lebesgue proofs via dominated
convergence of difference quotients; it does not silently equate
conditional improper integration with Lebesgue integrability.

\subsection{Reading formulas across
versions}\label{reading-formulas-across-versions}

Use the source IDs above, not a theorem number alone: Simon's 2014 and
2018 drafts number area and coarea differently. The older
\href{https://web.stanford.edu/class/math285/ts-gmt.pdf}{2014 course
draft}, printed p.54, formula 3.4, also has the incorrect target
dimension H\^{}m; the following weighted formula 3.5 has H\^{}n.~The
library cites the corrected 2018 area statement.

The convention in this directory is scalar integration against positive
measures. A parametrization's area Jacobian is nonnegative and counts
repeated coverings. An oriented vector flux additionally needs a chosen
normal and a signed dot product; it is not determined by the scalar
Hausdorff measure alone.

No source access, symbolic calculation, numerical test, or review
recorded here is a claim of proof-assistant verification.

\section{Verification ledger}\label{verification}

The collection contains 16 original worked expositions. Every example
received an independent mathematical review by a model agent other than
its author. These are recorded mathematical audits, with source and
computation checks; \textbf{no human peer review or proof-assistant
verification is claimed}.

\subsection{Independent review of each
example}\label{independent-review-of-each-example}

\begin{longtable}[]{@{}
  >{\raggedright\arraybackslash}p{(\columnwidth - 6\tabcolsep) * \real{0.1800}}
  >{\raggedright\arraybackslash}p{(\columnwidth - 6\tabcolsep) * \real{0.2200}}
  >{\raggedright\arraybackslash}p{(\columnwidth - 6\tabcolsep) * \real{0.2200}}
  >{\raggedright\arraybackslash}p{(\columnwidth - 6\tabcolsep) * \real{0.3800}}@{}}
\toprule\noalign{}
\begin{minipage}[b]{\linewidth}\raggedright
Examples
\end{minipage} & \begin{minipage}[b]{\linewidth}\raggedright
Author role
\end{minipage} & \begin{minipage}[b]{\linewidth}\raggedright
Independent reviewer role
\end{minipage} & \begin{minipage}[b]{\linewidth}\raggedright
Evidence and disposition
\end{minipage} \\
\midrule\noalign{}
\endhead
\bottomrule\noalign{}
\endlastfoot
E01--E04 & Lebesgue examples author & Geometry examples author &
\href{https://github.com/Drsunflowerherbert/Codex-for-Math/tree/codex/integral-examples-2026-10-15/integral-examples/reviews/review-E01-E04.md}{Review}:
all pass; E04's symmetric truncation is explicitly a subset of its open
domain. \\
E05, E06, E14--E16 & Parameters/\allowbreak{}convergence author & Area/coarea author
&
\href{https://github.com/Drsunflowerherbert/Codex-for-Math/tree/codex/integral-examples-2026-10-15/integral-examples/reviews/review-E05-E06-E14-E16.md}{Review}:
all pass; E14's real-measurable assumptions, diagonal subsequence, and
source references clarified. \\
E07--E10 & Geometry examples author & Parameters/\allowbreak{}convergence author &
\href{https://github.com/Drsunflowerherbert/Codex-for-Math/tree/codex/integral-examples-2026-10-15/integral-examples/reviews/review-E07-E10.md}{Review}:
all pass; no mathematical changes requested. \\
E11--E13 & Area/coarea author & Lebesgue examples author &
\href{https://github.com/Drsunflowerherbert/Codex-for-Math/tree/codex/integral-examples-2026-10-15/integral-examples/reviews/review-E11-E13.md}{Review}:
all pass; E13's empty-interior argument and description of shrinking
sets clarified. \\
\end{longtable}

The reviewer checked analytic assumptions and derivations, not just
final answers. The author signed off the recorded changes. Each review
contains SHA256 hashes of the exact reviewed Markdown files; the
\href{https://github.com/Drsunflowerherbert/Codex-for-Math/tree/codex/integral-examples-2026-10-15/integral-examples/reviews/review-manifest.json}{review
manifest} and collection checker detect changes after review. No
reviewer reviewed their own examples.

The
\href{https://github.com/Drsunflowerherbert/Codex-for-Math/tree/codex/integral-examples-2026-10-15/integral-examples/reviews/global-review.md}{global
review} separately audits the guide, Hausdorff normalization, theorem
dimensions and hypotheses, scalar versus oriented integration, source
locators, counterexample atlas, and retrieval metadata. One repaired
global statement now restricts sigma-finiteness to the Lebesgue/counting
product factors: ambient lower-dimensional Hausdorff measure need not be
sigma-finite. The source audit found dimensional misprints in older
notes and selected the corrected Simon 2018 area statement.

\subsection{Concrete computation
checks}\label{concrete-computation-checks}

Run \texttt{python3\ integral-examples/checks/check\_examples.py} from
the repository root. The recorded run in
\href{https://github.com/Drsunflowerherbert/Codex-for-Math/tree/codex/integral-examples-2026-10-15/integral-examples/checks/check-results.json}{check-results.json}
performed \textbf{92 concrete checks with zero failures}, using SymPy
1.14.0 and mpmath 1.3.0. The JSON records every identity, method,
expected value, result, input hash, package version, and computational
limitation.

\begin{longtable}[]{@{}
  >{\raggedright\arraybackslash}p{(\columnwidth - 4\tabcolsep) * \real{0.1200}}
  >{\raggedleft\arraybackslash}p{(\columnwidth - 4\tabcolsep) * \real{0.1400}}
  >{\raggedright\arraybackslash}p{(\columnwidth - 4\tabcolsep) * \real{0.7400}}@{}}
\toprule\noalign{}
\begin{minipage}[b]{\linewidth}\raggedright
Example
\end{minipage} & \begin{minipage}[b]{\linewidth}\raggedleft
Concrete checks
\end{minipage} & \begin{minipage}[b]{\linewidth}\raggedright
What those checks support
\end{minipage} \\
\midrule\noalign{}
\endhead
\bottomrule\noalign{}
\endlastfoot
E01 & 4 & Antiderivative, final value, series coefficient, quadrature in
the original order \\
E02 & 4 & Finite-stage length recurrence/sums and weighted product
computation \\
E03 & 0 & Proof-only: a nonmeasurable set and the null-envelope argument
are not computationally verified \\
E04 & 8 & Primitive identities, iterated values, finite rectangle bound,
section quadratures \\
E05 & 3 & Exact finite support sums and rectangular exhaustion
identities \\
E06 & 6 & Parameter derivatives, exponential integral,
positive/reversed/equal parameter quadratures \\
E07 & 5 & Speed, weighted length, unweighted length, quadratures \\
E08 & 5 & Speed/circulation density and once/twice trigonometric
integrals \\
E09 & 5 & Gram determinant, height and area integrals, quadratures \\
E10 & 5 & Spherical Gram determinant, moment, area, cap formulas \\
E11 & 6 & Fold determinant, multiplicity-weighted and once-counted
integrals \\
E12 & 11 & Radial Jacobian integrals, power thresholds/formulas,
selected quadratures \\
E13 & 8 & Exact finite-stage Cantor lengths/slopes/masses; no
infinite-stage theorem inferred \\
E14 & 4 & Exact mass and finite tail/convergence-set identities \\
E15 & 10 & Parameter derivative, antiderivative, symbolic endpoint
expression, scaled quadratures \\
E16 & 8 & Kernel mass, first moment, exponential test-function values,
quadratures \\
\end{longtable}

The mpmath checks use 70 decimal digits and an absolute-or-relative
comparison tolerance of \(10^{-50}\). This is a numerical consistency
test, \textbf{not a certified quadrature error bound}. Exact symbolic
checks are recorded separately.

Neither symbolic integration nor finite numerical sampling proves
measurability, absolute integrability, a limit exchange, a general fiber
count, or an infinite-stage Cantor assertion. Those are justified by the
written arguments and theorem hypotheses. E03 is explicitly marked
proof-only; E13's checks concern finite stages only. No Lean files are
supplied or Lean verification asserted for this collection.

\subsection{Structural and document
checks}\label{structural-and-document-checks}

\begin{itemize}
\tightlist
\item
  \href{https://github.com/Drsunflowerherbert/Codex-for-Math/tree/codex/integral-examples-2026-10-15/integral-examples/checks/check_collection.py}{check\_collection.py}
  validates all 16 IDs, mandatory sections, source IDs, metadata arrays,
  local file links, exact review hashes, computation-result hashes,
  retrieval commands, required deliverables, and the restriction to
  \texttt{integral-\allowbreak{}examples/\allowbreak{}} changes.
\item
  \href{https://github.com/Drsunflowerherbert/Codex-for-Math/tree/codex/integral-examples-2026-10-15/integral-examples/metadata.schema.json}{metadata.schema.json}
  documents the retrieval schema. The standard-library checker validates
  the required invariants directly; it does not claim to implement a
  complete JSON Schema engine.
\item
  \href{https://github.com/Drsunflowerherbert/Codex-for-Math/tree/codex/integral-examples-2026-10-15/integral-examples/build.py}{build.py}
  generates the retained LaTeX from the Markdown and compiles that exact
  LaTeX with pdfLaTeX. The
  \href{https://github.com/Drsunflowerherbert/Codex-for-Math/tree/codex/integral-examples-2026-10-15/integral-examples/checks/pdf-build.json}{build
  record} identifies source/output hashes, tool versions, and any
  remaining layout warnings.
\item
  \href{https://github.com/Drsunflowerherbert/Codex-for-Math/tree/codex/integral-examples-2026-10-15/integral-examples/checks/collection-check.json}{collection-check.json}
  records structural verification and final PDF/LaTeX hashes. PDF text
  and representative pages are also inspected during publication
  preparation.
\end{itemize}

\subsection{Publication scope and remaining
limits}\label{publication-scope-and-remaining-limits}

The base and fetched public branch heads are recorded in
\href{https://github.com/Drsunflowerherbert/Codex-for-Math/tree/codex/integral-examples-2026-10-15/integral-examples/checks/repository-base.json}{repository-base.json}.
The publication branch is \texttt{codex/\allowbreak{}integral-\allowbreak{}examples-\allowbreak{}2026-\allowbreak{}10-\allowbreak{}15},
based on \texttt{e4f337c3\allowbreak{}aa52ec4c\allowbreak{}3f2aab07\allowbreak{}8d7c91b3\allowbreak{}2ee0015f}. All
additions are under \texttt{integral-\allowbreak{}examples/\allowbreak{}}. Publication uses an
ordinary push to the new branch, without a pull request, force push,
merge, or modifications to other branches. The final task response
identifies the actual published commit and verification of remote
contents. This source document deliberately does not predict its own
future commit hash.

The earlier paper was not supplied or used. The collection is a checked
first batch, not an exhaustive atlas of all integration theorems. Its
general Lebesgue, area, and coarea theorems remain prerequisites, cited
in \hyperref[bibliography]{SOURCES.md}. Source PDFs are linked rather
than redistributed; online versions can change. Recorded review does not
remove the value of checking a theorem's precise hypotheses when
adapting an example to a new map, measure, signed weight, or unbounded
domain.

\end{document}
